
\chapter{Introduction}\label{chap:introduction}

The fifth generation of cellular communication, known as 5G, introduces substantial advancements over its predecessor, \gls{LTE}.
In contrast to the monolithic architecture of \gls{LTE}, 5G is built on a modular, microservice-based architecture that enables greater
scalability, flexibility, and ease of virtualization.
5G promises improvements in latencies, throughput, and energy efficiency.
In addition, a broader range of frequency bands is supported, from sub-\SI{6}{\giga\hertz} to millimeter-wave frequencies,
allowing for both wide-area coverage and high-speed transmissions.

Additionally, new use cases are defined for 5G communication, encompassing \gls{eMBB}, with a focus on
bandwidth and coverage, \gls{URLLC} for time-sensitive applications, and \gls{mMTC} for high device density and \gls{IoT} deployments.
The increased configurability compared to \gls{LTE} allows network operators to fine-tune parameters and tailor
the deployment directly to the service requirements.
For example, while the \gls{SCS} in \gls{LTE} is hard-coded, 5G introduces it as a configurable parameter.
This change in architecture directly influences system behavior and performance.

Another key introduction in 5G is the support for private, stand-alone deployments, so called campus networks.
The ability to operate private small-scale 5G deployments is a crucial innovation for industrial applications and the research community.
In manufacturing and logistics, for example, private 5G networks can enable reliable,
real-time communication between robotic systems or automated machinery, and thereby are a crucial building block for the Industry 4.0.

However, these advances also introduce complexity.
The flexibility and configurability of 5G systems create a significant challenge for researchers aiming to build
reproducible and consistent testbed setups and obtain comparable results.
Creating robust methods for 5G experiments is therefore essential to support the ongoing development of the technology.

Existing measurement studies tends to examine and control only a small range of parameters.
Additionally, the current literature frequently lack comprehensive descriptions of the testbed, including hardware and software versions, and configuration details.
This hinders the reproducibility and confirmability of the provided results.

One particular paper by Raffeck et al.\cite{raffeck2024parameterizing} is a key motivator for the research in this thesis.
In their work, the authors thoroughly documented their testbed setup and configurations.
Nevertheless, when comparing two testbeds with nearly identical hardware and the same software and configurations,
they observed significant performance differences.
This finding emphasizes the complexity and variability inherent in 5G systems.
Further it underlines the importance of a careful and fine-grained approach to 5G measurements.

To address these challenges, in this work, an automated framework for deployment, configuration and measurement of 5G campus networks is created.
The framework enables a systematic investigation of key parameters that affect reproducibility in 5G.
By covering a broad set of parameters, the influence of each parameter on system behavior can be analyzed in detail.
The automation of the deployment and the measurements ensures consistent system definitions and reproducible procedures.
Further, it enables efficient handling of large-scale measurement experiments.

The remainder of this work is structured as follows.
Chapter~\ref{chap:background} provides background information on 5G technology, the context to the collected  channel metrics and the definitions of reproducibility used throughout this work.
In Chapter~\ref{chap:relatedwork}, a selection of relevant research is presented and this thesis is positioned in the context of the existing literature.
Then, in Chapter~\ref{chap:methodology}, the methodology of the 5G deployment and automation is covered.
The obtained results are evaluted in two parts in Chapter~\ref{chap:evaluation}.
Section~\ref{sec:eval:validation} focuses on the validation of the system and presents results that analyze the impact of individual parameters on system behavior.
In contrast, Section~\ref{sec:eval:case} covers the outcomes of a factorial experiment that explores a broader set of parameters.
In Chapter~\ref{chap:future}, issues and limitations of the approach are discussed and possible avenues for future work are presented.
Finally, conclusions are drawn in Chapter~\ref{chap:conclusion}.


\chapter{Background}\label{chap:background}

In this chapter, the relevant background information required to understand this work, is covered.
First, in Section \ref{sec:back:5g}, a small overview over the functionality of 5G is given.
Then, in \ref{sec:back:channel}, channel metrics which are utilized in this work, are presented.
Finally, in \ref{sec:back:replicability}, the terms repeatability, reproducibility and replicability are defined.

\section{Fifth Generation of Cellular Networks}\label{sec:back:5g}

5G is the common name of the fifth generation of cellular mobile networks, defined by 3GPP.
It is fully specified since 2019 and brings significant changes from its predecessor \gls{LTE}.
While LTE was a monolithic design primarily geared towards mobile broadband, 5G introduces new areas of application,
as well as a shift to a service-oriented architecture.
These advancements result in a more flexible system, which allows operators to set the focus on different \gls{QoS}
aspects and enables 5G to support emerging applications.

To understand the remainder of this work, key aspects of the functionality of 5G are explained in this section.


\subsection{Use Cases}

A key trait of 5G is the relatively broad range of different applications.
There are three key scenarios of use for 5G, as defined by the International Telecommunication Union.\cites{series2017minimum, sachs20185g, zaidi2018ofdm}
\textbf{\Gls{eMBB}} is the advanced human-centric communication for mobile broadband, with the associated requirements of
high throughput and coverage.
\textbf{\Gls{mMTC}}, a set of technologies focussed on the support of \gls{IoT} devices.
For this use case, the focus is on the communication between a massive number of low-cost and battery-opperated devices.
Therefore, for this area of application, the focus is on the ability to provide service for a large
amount of devices, at low latency, and have high energy efficiency on the \gls{UE}.
\textbf{\Gls{URLLC}} is geared towards \gls{TSN}, requiring hard latency guarantees and high reliability for mission-critical applications,
like autonomous driving or factory automation.

To support these range of applications, the radio access technology for 5G, 5G \gls{NR}, features a high degree of configuration.
This allows the opperator to fit the system behavior to the requirements.
A system mutable to such a high degree provides researchers with a challenging task of achieving reproducibility and explainability of results.


\subsection{5G Core and Radio Access Network}

\begin{figure}[htb]
\centering
\includegraphics[width=0.95\textwidth]{/home/lks/DocSync/Uni/5G-Masterarbeit-Ausarbeitung/Ausarbeitung/res/figures/5G-Masterarbeit - overview.jpg}
\caption{Overview of the 5G system.}
\label{fig:back:overview}
\end{figure}


In 5G, the same high level components as in \gls{LTE} are used.
A \gls{UE}, consisting of a mobile station and a USIM, is the end device, i.e. a mobile phone.
The \gls{UE} is connected to the \gls{RAN} which in itself is connected to the core network.
The main entity of the \gls{RAN} is the radio transmitter, named \gls{gNB}.
The \gls{gNB} is connected to the 5G core.
Since this work performs measurements on the air interface, the functionality of the core is secondary and only briefly covered.
The \gls{5GC} itself is comprised of separate network functions, which can be virtualized and scaled up or down independently of another.
Further, the \gls{5GS} can run in \gls{NSA} mode, where the 5G \gls{RAN} is used in conjunction with existing
\gls{LTE} infrastructure, however, the full feature set of 5G is only available in \gls{SA} mode.\cite{3gpp5goverview,3gpp23.501}
In this work, however, only a \gls{SA} architecture is considered.

Traditionally, a \gls{RAN} is a monolithic solution, combining a \gls{RU} and a baseband unit.
It is provided by vendor and operates on proprietary interfaces and standards.
It can be highly optimized, however due to the monolithic nature, it is inflexible.
The Open RAN approach aims to open the \gls{RAN} and provide a set of standards,
which allows the interoperability of different building blocks of the \gls{RAN}.
Splitting up the \gls{RAN} into multiple well defined modules provides a higher degree of flexibility, cost reduction
and an increase in competition.
In Open RAN, the baseband unit is split into a central unit and a distributed unit.
Further the protocols between the \gls{5GC} and the central unit, as well as the protocol between the central unit and
the distributed unit and the protocol between the distributed unit and the \gls{RU} are standardized, enabling the
interoperability between different vendors.\cite{polese2023understandingopenran}
For this work hoever, a traditional RAN design is used.

\subsection{5G New Radio}

5G \gls{NR} is the standard for the high performance air interface of the fifth generation of cellular communication systems.
Since all measurements in this work are performed on the air interface, understanding its structure and behavior is essential.
This section introduces the key elements of 5G \gls{NR} that are relevant to the measurement process.

It begins with the composition of the 5G frame structure in Section~\ref{sec:back:frame}, followed by an explanation
of the 5G numerology and its effect on the subframe division in Section~\ref{sec:back:numerology}.
Then, in Section~\ref{sec:back:duplexing}, the duplexing schemes, which coordinate the uplink and downlink transmissions, are presented.
Finally, in Section~\ref{sec:back:resource} the structure of the ressource grid is introduced, which explains how time and frequency resources are allocated in 5G.


\subsubsection{Frame Structure}\label{sec:back:frame}
The frame structure of 5G \gls{NR} is an extension from \gls{LTE} and is designed for flexibility.
In Figure \ref{fig:back:tddpattern}, the frame structure can be seen.
The highlighted duplexing pattern will be explained in Section \ref{sec:back:duplexing}.
The largest scheduling unit is a frame with a duration of \SI{10}{\milli\second}.
Each frame consists of 10 subframes, each with a duration of \SI{1}{\milli\second}.
Depending on the \gls{SCS}, each subframe is further divided into slots.
In 5G, as opposed to \gls{LTE}, the number of slots per subframe can be configured and depends on the \gls{SCS}.
The exact mechanisms is explained in Section \ref{sec:back:numerology}.
Each slot further contains 14 symbols, or 12 symbols for an extended cyclic prefix with a longer guard interval.
Depending on the type of duplexing used, symbols either are all used for uplink or downlink, regulated by the carrier,
as is the case for \gls{FDD}.
Alternatively, for \gls{TDD}, each \gls{OFDM} symbol can be used for uplink or downlink, however guard symbols are
required when changing from downlink to uplink.
\cite{hui2018channel,3gpp38.211,3gpp38.912}
% The different duplexing mechanisms in 5G are explained in more detail in Section \ref{sec:back:duplexing}

\begin{figure}[htb]
\centering
\includegraphics[width=0.95\textwidth]
{/home/lks/DocSync/Uni/5G-Masterarbeit-Ausarbeitung/Ausarbeitung/res/figures/5G-Masterarbeit - TDD pattern_2.jpg}
\caption{
    Representation of 5G frame structure ($\mu = 1$) with color coded TDD pattern: \setulcolor{Cgreen1}\ul{downlink (green)},
\setulcolor{Cplotn1} \ul{uplink (orange)}, \setulcolor{Cblueish3}\ul{flex slot (blue)},
\setulcolor{Cpinknew0} and \ul{guard symbol (pink)}.
}
\label{fig:back:tddpattern}
\end{figure}


\subsubsection{Numerology}\label{sec:back:numerology}

A significant advancement from \gls{LTE} towards 5G \gls{NR} is the change of hard-coded parameters to configurable.
Thereby enabling the operators to dynamically fit the configuration to the current requirements.

One example of this flexibility can be seen in the 5G \gls{NR} numberology, in particular the options for \gls{SCS}.
As in \gls{LTE}, frames have a fixed duration of \SI{10}{\milli\second}, each consisting of 10 subframes with a duration of \SI{1}{\milli\second}.
In contrast to \gls{LTE}, the amount of slots per subframe can be adjusted.
The parameter $\mu$ is used to denote the \gls{OFDM} \gls{SCS} with a frequency of
$ \SI{15}{\kilo\hertz} \cdot 2^{\mu}$.
At the time of this work, $\mu$ is defined for values of 0 to 6.
The amount of slots per subframe and therefore the duration of each slot and symbol depends on the selected
\gls{SCS}, with the amount of slots per subframe being equal to $ 2^\mu $. \cite{3gpp38.211, zaidi2018ofdm}
In this work, only the effect of the \gls{SCS} on the slot duration is relevant, therefore additional details omitted.

\subsubsection{Duplexing}\label{sec:back:duplexing}

In 5G \gls{NR} there are two posibilities to duplex uplink and downlink traffic, \gls{FDD} and \gls{TDD}.
For \gls{FDD}, distinct frequency bands for uplink and downlink transmissions are used, providing full-duplex.
This however comes at the cost of a lowered spectral efficiency, since a guard band between the uplink and
downlink carrier is required.
For \gls{TDD}, a period length and a \gls{TDD} pattern must be defined.
The period length describes a number of consecutive slots and must result in a duration of \SI{0.5}{\milli\second},
\SI{0.625}{\milli\second}, \SI{1.25}{\milli\second}, \SI{2.5}{\milli\second},
\SI{5}{\milli\second}, or \SI{10}{\milli\second}.
For a \gls{SCS} of $\mu = 1$ and a resulting slot duration of \SI{0.5}{\milli\second}, valid period lengths are
therefore 5, 10, and 20 slots.
The \gls{TDD} pattern describes the allocation of slots to downlink, uplink, or flexible.
Furthermore, for flexible slots, the symbols in a slot do not have to belong to the same type.
Symbols in a slot can be assigned to downlink, uplink or flexible.
Flexible symbols can be used dynamically for downlink, uplink, or as guard symbol to signal the switch from
downlink to uplink transmission.\cite{3gpp38.213,3gpp38.912}
An example \gls{TDD} configuration for $\mu = 1$ and a period length of 10 slots can be seen in Figure \ref{fig:back:tddpattern}.
It displays how a 5G \gls{NR} frame is split into 10 subframes, which in turn split into 20 slots.
The first five slots are allocated for downlink traffic, then a flex slot, followed by four uplink slots.
This \gls{TDD} pattern with a length of 10 slots is then repeated.
The flex slots, slot five and 15, contain three downlink symbols, then two guard symbols, to signal the switch from
downlink to uplink, followed by nine uplink symbols.

With an increasing distance between \gls{gNB} and \gls{UE}, longer guard times are required,
making \gls{FDD} a better option for wide coverage.\cite{berardinelli2016guardtimetdd}
\Gls{TDD} on the other hand offers better flexibility and spectral efficiency, therefore it is well suited
for campus networks and deployments with a limited range.
For the n78 band used throughout this work, only \gls{TDD} is defined, therefore no \gls{FDD} measurements are
performed.\cite[.p~44]{3gpp38.104}

\subsubsection{Resource Grid}\label{sec:back:resource}

\begin{figure}[htb]
\centering
\includegraphics[width=0.95\textwidth]{/home/lks/DocSync/Uni/5G-Masterarbeit-Ausarbeitung/Ausarbeitung/res/figures/5G-Masterarbeit - resource grid_r.jpg}
\caption{Resource Grid in 5G \gls{NR}.}
\label{fig:back:resourcegrid}
\end{figure}

The resource grid in 5G \gls{NR} is mostly similar to the resource grid in \gls{LTE},
however the actual scale of the dimensions can vary, depending on the numerology.

The 5G resource grid spans two dimensions, the frequency dimension and the time dimension.
One excerpt of an example resource grid can be seen in Figure \ref{fig:back:resourcegrid}.
In the time dimension, there are frames, consisting of subframes which consists of slots which in itself typically
contains 14 symbols, as described in Section \ref{sec:back:frame}.
The 14 symbols of one slot can be seen along the vertical axis in the figure, the frequency is shown along the horizontal axis.

In the frequency dimension, the largest unit is a \gls{PRB}.
A \gls{PRB} is made up of exactly 12 consecutive subcarriers and therefore occupies a bandwidth
of $ 12 * \gls{SCS} $, i.e. $ 12 * \SI{30}{\kilo\hertz} = \SI{360}{\kilo\hertz}$ for $ \mu = 1 $.
Dependent on the available and desired bandwidth, multiple \gls{PRB} are defined back-to-back.
The smallest unit on the resource grid is called a \gls{RE} and is made up of one subcarrier of one \gls{PRB} in
the frequency domain and one OFDM symbol in the time domain.\cite{3gpp38.912,3gpp38.211}

For each antenna port, \gls{SCS} configuration, and transmissions direction a resource grid is defined.


\subsection{5G Campus Networks}

In this work, the focus is on private 5G \gls{SA} campus networks.
The ability to operate private, stand-alone 5G networks is a key feature of 5G.
A private 5G network is a local area network, which utilizes 5G \gls{NR} technology within a defined geographic region (campus).\cite{aijaz2020private}
By utilizing a private 5G network, the operators have immediate control over the network.
Without interference from undesired devices, the network can be taylored to the required service quality like latency, reliability and availability.
Therefore, this makes 5G campus networks attractive for industrial applications in various sectors.


\section{MCS, SNR, RSRP and other Channel Quality Indicators}\label{sec:back:channel}

In 5G, there are multiple \gls{CSI} defined, which can reported by the \gls{UE} to the \gls{gNB}.
These \gls{CSI} may include \gls{CQI}, \gls{PMI}, CSI-RS Resource Indicator,
SS/PBCH Block Resource Indicator, Layer Indicator, \gls{RI}, Layer 1 \gls{RSRP}, or Layer 1 \gls{SiNR}.\cite{3gpp38.214}
Information from the \gls{CSI} report is used for decisions affecting the link modulation, beam management and resource scheduling.
Since the \gls{gNB}s in this work only utilizes \gls{SISO}, as opposed to \gls{MIMO}, where multiple antennas are used simultaneously,
most of the \gls{CSI} metrics are secondary.
Relevant for this work are \gls{RSRP}/\gls{RSRQ}, \gls{SiNR}/\gls{SNR}, \gls{CQI}, as well as \gls{MCS}, which
encapsulates the selected modulation and channel coding, based on the received \gls{CSI}.

\subsection{Modulation Coding Scheme}\label{sec:back:channelmcs}

The \gls{MCS} is determined based on the \gls{CQI} and corresponds directly to a modulation scheme and a code rate.
There are multiple lookup tables defined, based on the maximum supported modulation scheme, i.e. 64\gls{QAM} or 256\gls{QAM}.
Dynamically reacting to changes on the air interface allows 5G \gls{NR} to minimize \gls{BLER} and retransmissions
while simultaneously optimizing throughput.
A higher \gls{MCS} index results in a more efficient modulation, from \gls{QPSK} to 16\gls{QAM},64\gls{QAM}, and 256\gls{QAM}.
Within the same modulation scheme, a higher \gls{MCS} index also corresponds to a higher code rate with less redundancy bits.\cite{dahlman5GNR,3gpp38.214}

\subsection{Signal to Noise Ratio and Signal to Interference plus Noise Ratio}

\Gls{SNR} and \gls{SiNR} are both general metrics of channel quality.
The \gls{SNR} describes the ratio between the signal and noise power, while the \gls{SiNR} describes the ratio of signal
power to the sum of interference power from other cells and noise power.
Higher values are better and are typically provided in the logarithmic measure \SI{}{\deci\bel}.
A \gls{SNR} ratio of \SI{20}{\deci\bel}, for example, states a ratio of 100:1.
The \gls{SiNR} is measured on the \gls{UE} on a \gls{PRB} basis. Based on the computed \gls{SiNR} the
\gls{CQI} is calculated.\cite{afroz2015sinr}


\subsection{Reference Signal Received Power and Reference Signal Received Quality}

In cellular networks, \gls{RSRP} and \gls{RSRQ} report the received power and the quality of the received reference
signal, respectively.
\gls{RSRP} indicates a good signal strength in the range of \SI{-90}{\dBm} to \SI{-80}{\dBm} and
good signal quality reports \gls{RSRQ} values of \SI{-15}{\deci\bel} to \SI{-10}{\deci\bel}.
In both cases, higher values are better.
\gls{RSRP} and \gls{RSRQ} are used in decision making for cell selection and handover.
Both metrics are cell-specific, while typically, \gls{RSRP} is more significant and \gls{RSRQ} mainly
employed if \gls{RSRP} measurements are not sufficient.\cite{lazar2023realchannel,afroz2015sinr}

\subsection{Channel Quality Indicator}

In 5G, \gls{CQI} is quality indicator for the communication channel.
Periodic decoding of the reference signal by the \gls{UE} is used to determine the \gls{CQI}.
As a result, the \gls{CQI} is a quantized and scaled measure of the experienced \gls{SiNR}.
\gls{CQI} ranges from 0 to 15 with a higher value indicating the better channel quality.
The \gls{gNB} uses the reported \gls{CQI} from the \gls{UE} to determine the adequate \gls{MCS}.\cite{afroz2015sinr}



\section{Repeatability, Reproducibility and Replicability}\label{sec:back:replicability}

The Association for Computing Machinery defines Repeatability, Reproducibility and Replicability as follows.\cite{acm_reproducibility}

Repeatability refers to the ability to obtain consistent results, within the stated precision, the same measurement
procedure, equipement, and operating conditions, at the same location, for multiple repetitions, performed by the same team.
This concept is often summarized as \emph{"Same team, same experimental setup"}.

Reproducibility denotes the degree to which identical results can be achieved, within the stated precision,
with the same measurement
procedure and system, under the same operating conditions, the same or a different location, for multiple repetitions, by a different team.
Reproducibility is commonly abbreviated to \emph{"Different team, same experimental setup"}

Replicability is defined as the ability to obtain the same measurement results, with the stated precision, a different
measureing system, a different location, for multiple runs, by a different team.
Therefore, an independent team is able to optain the same results independently.
Replicability is often abbreviated to \emph{"Different team, different experimental setup"}.



\chapter{Related Work}\label{chap:relatedwork}

A comprehensive understanding of the factors influencing 5G \gls{QoS} is not yet fully understood and remains an area requiring further investigation.
However, there is already a substantial body of related work regarding real-world measurements in 5G networks.

A modular and flexible testing methodology is presented in~\cite{diaz2020modular}.
The approach bundles all experimental components, such as test cases, scenarios, and slices, for automatic management and execution.
By applying this methodology to the 5G \gls{NSA} deployment at the University of Malaga,
crucial insights into the local setup were gained.
The work also highlights the overarching need for dependable testing and validation of 5G services.

The work in~\cite{foukas2018experience} describes a 5G testbed composed of an edge
cloud at the University of Edinburgh and a central cloud at King's College London.
The core network in this setup is powered by Orion,
a RAN slicing system built on top of OAI and developed by the authors.
The authors identify key areas for future research,
including observability and orchestration.

An investigation into the throughput and delay of 5G NR is conducted in~\cite{raffeck2024parameterizing}.
For this purpose, two separate teams performed identical measurements at different locations,
utilizing comparable testbeds.
The testbeds utilize the same \gls{gNB} and \gls{UHD} versions and are equipped with the same modem and SDR.
Nonetheless, the results revealed significant performance differences between the two testbeds in certain analyzed scenarios.

In~\cite{Rischke20215GCN}, an extensive measurement study of 5G \gls{QoS} is performed.
This research was motivated by the need for packet-level measurements in \gls{NSA} and \gls{SA} 5G deployments.
Packet mirroring at an intermediate switch is used to enable high precision packet captures.
The study is conducted with \gls{CBR} traffic evaluates \gls{OWD}, packet loss, \gls{IPDV}, and downtime.
The evaluation presents a general wide spread of delay, significant outliers in core network delay, and non-symmetrical upload and download delays.
The authors suggests that \gls{RAN} scheduling is not yet suitable for industrial applications and refers to future work for this topic.

The latency in 5G NR Standalone campus networks is analyzed in ~\cite{Rischke2022EmpiricalSO}.
The experimental setup consists of a Nokia ABIL+ASIK gNB and a Nokia AZQH antenna, while the core is the Open5GS SA implementation.
Throughout the experiments, different combinations of \gls{PPS}, packet sizes and traffic distributions are used.
In addition to the \gls{OWD} and the resulting \gls{RTT}, the \gls{IPDV} and density is evaluated.
A key conclusion of this work is the negative correlation observed between uplink and downlink latency, as well as a
discrete distribution of \gls{RTT} for \gls{CBR} traffic.

A performance evaluation of a 5G deployment utilizing srsRAN and open5GS is detailed in \cite{10439170}.
The hardware configuration includes an Ettus USRP B210 \gls{SDR} and two different modems as well as one mobile phone as \gls{UE}.
The evaluation analyzes network coverage, throughput, and \gls{RTT} latency.
To assess performance under different conditions, key parameters, like the \gls{TDD} pattern and modulation scheme, are modified.

In \cite{Zhao2023PhysicalLR}, a latency analysis of 5G NR, with a focus on ultra-low latency and the physical layer, is performed.
It provides a comprehensive overview of various delay types and their contributing factors.
Furthermore, the paper offers estimations regarding the specific contribution of each latency type to the overall delay.

A data-driven test platform designed for 5G research is introduced in \cite{Wang2021DevelopmentOA}.
It featuring a data collection process optimized for machine learning applications.
The paper illustrates the platform's utility through two use cases.
The first case aims to prevent unnecessary cell switches between LTE and NR in a \gls{NSA} deployment.
The second use case demonstrates antenna direction recognition, which enables beam-forming without requiring UE
location information.

Best practices and a troubleshooting guideline for deploying 5G Standalone networks are the covered in
\cite{Mamushiane2023DeployingAS}.
This guideline is focussed on setups using an Open5GS core in combination with srsRAN.
The paper offers practical recommendations for 5G deployments, covering aspects like compute node configuration,
5G UE integration, and network slicing.

An in-depth analysis of latency on the 5G mmWave physical layer is conducted in \cite{fezeu2023depth}.
The authors performed measurements using different commercial phones as UEs within a commercially operated 5G network.
The paper identifies several factors that substantially influence delay,
including channel conditions and UL scheduling mechanisms,
which hinder the goal of achieving reliable sub-millisecond latency.

In \cite{mamushiane2023towards}, the resource requirements of the Open5GS Core are analyzed.
In this setup, \gls{UE} and \gls{RAN} are emulated via the UERANSIM stack.
This work specifically focuses on benchmarking and investigating the CPU and RAM utilization of the \gls{UPF} when placed under stress.

A comparison of delay characteristics among different modems with varying chipsets is presented in \cite{seidel2024impact}.
The underlying testbed utilizes OAI for both, the \gls{RAN} and the core network.
The study confirms significant differences between 5G chipsets, affecting the baseline delay and packet distribution.
A particularly notable discovery was a distinct sawtooth pattern observed in the packet distribution of the Qualcomm X55 chipset.

The performance of the 5G SA campus network at the University of Kaiserslautern is analyzed in \cite{mallikarjun2022performance}.
This network's RAN and CN are both supplied by Nokia.
Using three distinct UEs, the study measured throughput, latency, and signal strength parameters including RSS, SS-RSRP, and SS-SINR.
The results revealed significant performance variations depending on the measurement point and the specific UE used.

In \cite{doria2024virtualized} a 5G testbed using OAI is evaluated.
For the setup, the authors provide detailed configuration descriptions.
Its performance are assessed across three different measurement scenarios: spectrum occupation,
coverage evaluation, and throughput and RTT measurements.
During these tests, parameters such as the number of \gls{PRB} and the utilized bandwidth are controlled.

An analysis of the resource consumption of the OAI 5G core is presented in \cite{atalay2022scaling}.
The experimental setup involves the OAI core network alongside the gNBSIM RAN simulator.
This research focusses on system behavior when utilizing network slicing and changing the system load.
A key part of the study is the recording of memory and CPU consumption for individual network functions

In \cite{kim2024toward}, an approach to reproduce field condition 5G MIMO OTA tests in a lab.
The approach is based on capturing, analyzing and reconstructing the 5G signals.
The paper confirms the general feasibility of recreating field test conditions in the lab.
However, future work and testing, e.g. with an anechoic chamber, is required to enable precise replication.


\bigskip

Among the presented papers where 5G testbed measurements are crucial,
several provide an in-depth description of the hardware and software in the setup~\cite{diaz2020modular, raffeck2024parameterizing, 10439170, Mamushiane2023DeployingAS, mallikarjun2022performance, doria2024virtualized}.
A smaller subset of these offers a thorough breakdown of the configuration aspects~\cite{diaz2020modular, raffeck2024parameterizing, 10439170, doria2024virtualized}.

In the related work, a range of heterogeneous 5G testbeds are presented.
The setups include all-in-one 5G boxes \cite{Wang2021DevelopmentOA},
proprietary hardware solutions from vendors like Nokia \cite{diaz2020modular, mallikarjun2022performance},
and hybrid approaches that combine proprietary systems and open-source software or accessible hardware \cite{Rischke20215GCN, Rischke2022EmpiricalSO}.
Furthermore, many studies utilize completly accessible testbeds, built with open-source software in combination with a \gls{SDR} \cite{foukas2018experience, raffeck2024parameterizing,10439170,Mamushiane2023DeployingAS,mamushiane2023towards,seidel2024impact,doria2024virtualized,atalay2022scaling}.

These aspects are crucial for enabling a third party to replicate research findings.
However, of all the covered work, only three papers \cite{raffeck2024parameterizing,10439170,doria2024virtualized}
provide all the elements for convenient reproducibility: an in-depth, detailed testbed description,
thorough coverage of configuration options, and the use of open-source and accessible hardware and software.

The full range of factors influencing 5G performance is not yet completely understood,
as much of the existing research controls for only a limited set of parameters.
{In addition, there are multiple guidelines covering QoS and security of 5G, however, to the knowledge of the authors,
no 5G guideline focussed on the reproducibility has been published so far.}
Key variables are often overlooked, including environmental factors like background radiation and proximity
to reflective surfaces, as well as system-level differences such as varying traffic patterns, packet sizes, and differences between software versions or hardware components.


The significance of this lack of explainability is highlighted by Raffeck et al. in~\cite{raffeck2024parameterizing}.
They demonstrate, that two teams can produce significantly different measurement results, even when using comparable hardware with identical software and configurations.
Although expanding the number of observed parameters increases the complexity of the work and the time required for measurements,
a detailed investigation into both environmental and system factors is essential for obtaining valid and replicable results.

Therefore, this work provides an in depth investigation into the reproducibility in 5G campus networks.
In detail, this work is a continuation of the research presented in \cite{raffeck2024parameterizing} and provides
\textbf{a)} an automated and straightforward repeatable setup suite for 5G campus deployments.
\textbf{b)} an automated measurement suite, capable of handling a wide range of parameters.
\textbf{c)} a repetition of the measurement results demonstrated in~\cite{raffeck2024parameterizing}.
\textbf{d)} an investigation into several key scenarios and the effect of certain parameters and parameter combinations.
\textbf{e)} a case study, covering a wide range of parameters and an evaluation of their interactions.


\chapter{Methodology}\label{chap:methodology}

This chapter provides the methodology employed in this work, structured into three main sections.
First, the hardware and software components of the testbed, as well as the associated configuration
is explained in Section~\ref{sec:meth:testbed}.
Then, the measurement procedure is illustrated in Section~\ref{sec:meth:measure}.
In this context, all measurement parameters are explained and the steps taken to create an automatable setup are highlighted.
Finally, in Section~\ref{sec:meth:issues}, known technical issues are listed and mitigations are presented.

\section{Testbed}\label{sec:meth:testbed}

\begin{figure}[htb]
\centering
\includegraphics[width=0.90\textwidth]{/home/lks/DocSync/Uni/5G-Masterarbeit-Ausarbeitung/Ausarbeitung/res/figures/IMG_8977_c01_ff01.jpg}
\caption{Picture of the actual testbed.}
\label{fig:meth:testbedPicture}
\end{figure}

In this work, a 5G campus network is used to investigate the influences of certain parameters.
This setup is presented in Figure~\ref{fig:meth:testbedPicture}.
On the tripod in the foreground, the \gls{SDR} can be seen, while the modem is visible on the tripod in the background.
A schematic representation of the setup can also be seen in Figure~\ref{fig:meth:testbedSetup}.
The testbed consists of two machines, connected via Ethernet.
One machine is running the Open5GS \gls{SA} \gls{5GC}.
The other machine operates as the \gls{gNB}.
An \gls{SDR} is attached via USB to the \gls{gNB} host and functions as the \gls{RU}.
Additionally attached to the \gls{gNB} host is an USB modem, which is the \gls{UE} in this setup.
This allows timestamping packets on the \gls{gNB} host with \texttt{tcpdump}.
Placing a USB hub between the modem and the \gls{gNB} host gives the ability to remotely powercycle the modem.
Measurements are performed at the interface which corresponds to the modem and at the interface connected to the \gls{5GC}.
Modem and \gls{SDR} are placed on tripods to enable control of height and distance and allow precise positioning.

\begin{figure}[htb]
\centering
\includegraphics[width=0.90\textwidth]{/home/lks/DocSync/Uni/5G-Masterarbeit-Ausarbeitung/Ausarbeitung/res/figures/testbed-setup.pdf}
\caption{Testbed setup.}
\label{fig:meth:testbedSetup}
\end{figure}


\subsection{Testbed configuration}

All experiments in this work are conducted within the n78 band at a frequency range between \SI{3700}{\mega\hertz}
and \SI{3800}{\mega\hertz}.
A channel bandwidth of \SI{20}{\mega\hertz} with a center frequency of \SI{3750}{\mega\hertz} is used.

To ensure consistent measurement conditions, both the \gls{RU} and the \gls{UE} are placed at a height of \SI{1}{\meter}.
Apart from one specific scenario, a minimum distance of \SI{1}{\meter} from walls and other surrounding objects
is maintained throughout all scenarios to reduce the influence of reflections and interference.

The \gls{RU} is operated in \gls{SISO} mode for all measurements.
Therefore, features specific to \gls{MIMO} configurations, like beamforming, are outside the scope of this work and were not considered.


\subsubsection{Antenna Gain}\label{sec:meth:antennagain}
The quality of the transmission channel and thereby the achieved offered throughput in part depends on the configuration
of the RX and TX gain settings of the \gls{RU}.
Suitable values can be determined experimentally and may be affected, among others, by the distance between \gls{RU}
and \gls{UE}, the type of antenna, distance to walls, floors, and other reflective surfaces, as well as background noise.
Therefore, the RX and TX gain values must be redetermined for each new location.

In this work, for each configuration which differs in location or distance, a new configuration measurement
was performed beforehand.
For this, one of the to parameters, RX gain and TX gain, was held at a reasonable default, while the other value
was varied in \SI{4}{\deci\bel} steps; the roles where then reversed.
In theory, suitable RX gain and TX gain values can not be determined independently of each other,
requiring a factorial experiment.
In this work, due to time constraints, the interactions between RX gain and TX gain were not considered.
In the future, an approach utilizing gradient descent could be used to determine a suitable RX/TX gain combination
at a reasonable time.

\subsubsection{SIM Card Provisioning}

For this setup, open SIM cards from the Open Cells Project~\cite{opencells} are used.

Flashing the SIM card by using osmocoms pySim~\cite{pysim} was not successfull and resulted in a corrupted card.
Instead \texttt{program\_uicc}, a small C program from the open cells project~\cite{opencellssw} was used.
The setting up of a new SIM card, with an \gls{IMSI} and a set of cryptography parameters can be seen in Listing~\ref{lst:meth:simflash}

\begin{lstfloat}[htb]
\begin{lstlisting}[basicstyle={\ttfamily\small\color{black}},style=base,caption={Flashing new IMSI to SIM.},label={lst:meth:simflash}]
@$@ ./program_uicc --adm 12345678 --imsi 999420000014671 \
    --key BDFDFD2BE954A1AA29765DB6DAEEF5E7 \
    --opc 358422278845A5632BBFB7B354DB103A \
    --spn lsinfo3
\end{lstlisting}
\end{lstfloat}


Matching KI/OPc pairs, the required cryptography parameters, can be generated with
\texttt{kiopcgen}~\cite{kiopcgenerator} as shown in Listing~\ref{lst:meth:simkey}.

\begin{lstfloat}[htb]
\begin{lstlisting}[basicstyle={\ttfamily\small\color{black}},style=base,caption={Generated KI/OPC pairs by usind python kiopcgen.},label={lst:meth:simkey}]
@$@ python kiopcgen --op D7DECB1F50404CC29ECBF989FE73AFC1 \
    --transport 2211CC6E9746434B89F346F0276CCAEC
{'KI': '8BB1FEC063CC436EA6C1BF2CAD450251',
 'OPC': b'A3F637D7F3D3D4CDB818A7B98BEDB9C6',
 'eKI': 'B9D6C1CE933359BE0A52BAC9E81CB6D1'}
\end{lstlisting}
\end{lstfloat}

The WebUI of the Open5GS \gls{5GC} can be used to add the newly provisioned SIM card as a subscriber.
This step requires the \gls{IMSI} and the Ki/OPc pair.


\subsection{Hardware and Software Specification}
The machines for the \gls{5GC} and the \gls{gNB} are both running Ubuntu 22.04 and are equipped with an
AMD Ryzen 5 PRO 5650G processor and \SI{16}{\giga\byte} of RAM.
The modem used in this work is an WaveShare SIM8200EA-M2 5G which is based on a Qualcomm X55 platform.
The \gls{RU} is an Ettus USRP B210 \gls{SDR}.
The USB hub enabling the power cycling of the USB modem is the MEGA4: 4-Port USB 3.1 PPPS Hub for Raspberry Pi 4B from UUGear.
Depending on the scenario, \gls{UHD} version 3.15 LTS and version 4.0 LTS are used for the \gls{SDR}.
For OAI, version 2.1.0 (commit \texttt{9fab212}) and version 2.2.0 (commit \texttt{6819108}) and for
srsRAN versions 24.04 and 24.10 are deployed.
The \gls{5GC} is running the Open5GS \gls{SA} application release 2.7.2.


\section{Measurement Setup}\label{sec:meth:measure}

In this work, automated measurements, based on a definition of all required parameters, are performed.
For each measurement run, \SI{60}{\second} of traffic are recorded.
The long measurement time makes the recorded results less susceptible to one-off events and flucuating
behavior in early stages of the connection.
The total time required for a measurement is approximately \SI{3.5}{\minute}, mainly dominated by the modem's
reinitialization time after a power cycle.
As a result, reducing the measurement duration has limited impact on the total number of measurements
that can be performed within a fixed timeframe.

A minimum of five repetitions if performed for each observed parameter combination.


\subsection{Ansible and Measurement Automation}\label{sec:meth:measureautomation}
To produce a clean and reproducible measurement procedure, automation is implemented via Ansible playbooks.\cite{ansible}
Ansible is a widely adopted automation system, commonly tasked with configuration management,
provisioning, and deployment.
Ansible provides modules for most common tasks, like package management, handling of systemd services, and text manipulation.
Collecting task in Ansible playbooks leads to a self documenting workflow with good error handling abilities.
The error handling is provided at a per module basis and in general, which easily enables gracefull recovery.
Ansible playbooks are written in yaml, which is human-friendly and simultaneously well parsable by machines.
Additionally, Ansible is built to provide declarative definitions instead of procedural.
This makes Ansible a good fit for automation in a context where error-resistance and reproducibility are required.

Each measurement run is completly defined by a set of system and configuration parameters.
Multiple measurement runs can be combined into a series, as long as no parameter requires physical intervention,
e.g. by changing the distance between \gls{RU} and \gls{UE}.
Thus, a series is a combination of fixed system parameters and a list of parameters specific to each run.
The format of this definition can be seen in Listing~\ref{tab:meth:definitionmsm}.
Additionally, an excerpt of the per-run definition can be seen in Listing~\ref{tab:meth:definitionmsmSingleRun}.
To group and identify measurement results, the set of system parameters is hashed and the resulting hash
is used as an identifier.
The identifier for each run is a combination of the system identifier, the hash of the per-run parameters,
and the number of repetitions for this parameter combination.

A measurement series is the basis for the automation of measurement runs.
The Ansible playbook is given a definition of a measurement series and loops over all runs.
After each run is completed, the run ID and a timestamp is written to a log file.
To compensate the issues described in Section~\ref{sec:meth:issues}, particularly the Issue described
in Section~\ref{sec:meth:issueinterfacename}, a small bash script is used to execute the Ansible playbook.
This bash file catches unclean exits of Ansible, reboots the \gls{gNB} host, parses the log to get
the ID of the last completed run, and triggers the continuation of the Ansible playbook.



\subsection{Measurement Run}\label{sec:meth:measuresingle}
In contrast to the high level insight into the measurement automation provided in Section~\ref{sec:meth:measureautomation},
this section explains the concrete steps required to perform a singular measurement run.

\begin{figure}[htb]
\centering
\includegraphics[width=0.95\textwidth]{/home/lks/DocSync/Uni/5G-Masterarbeit-Ausarbeitung/Ausarbeitung/res/figures/5G-Masterarbeit - Per run-2.pdf}
\caption{Logic of programm flow for a single measurement run.}
\label{fig:meth:automationrun}
\end{figure}

In Figure~\ref{fig:meth:automationrun} a schematic of the process can be seen.
For reasons of clarity and comprehensibility, some details are omitted and subroutines,
consisting of multiple tasks are combined.
Subroutines are marked by the dashed border.
Additionally, steps which are immediately dependent on the parameters for this run are highlighted with a colored border.

After the definition of a measurement series is parsed, the Ansible playbook loops over the defined runs.
As a precaution, it is verified, that all applications from the previous run, like \gls{gNB} or logging scripts, are terminated.
Then, the correct version of \gls{UHD} is loaded.
To recreate the same state every time, the modem is power cycled before the run.
There are some caveats to this procedure, which are explained in Section~\ref{sec:meth:modemcycle} in more detail.
Then, the configuration of the selected \gls{gNB} type and version is modified, depending on the configured \gls{TDD} pattern and RX/TX gain values.
The configuration is copied to the target machine and the \gls{gNB} application is started with the new configuration.
Alternatively, the wanted docker container is started.
After the \gls{gNB} has been started, ModemManager is used to connect the modem to the \gls{5GC} and acquire an IP address.
If the \gls{gNB} exits due to an error or the modem can't acquire an IP address, the run is marked as failed.
The reason for the failure is saved and temporary files are cleaned up.
If the modem has successfully connected, a script which periodically polls channel metrics from the modem is copied
to the target machine and then started in the background.
The method used to acquire modem channel metrics is described in Section~\ref{sec:meth:channelmetrics}.
Then the routing informations are set.
The executed traffic role depends on the type of traffic, i.e. iperf3 for throughput or scapy for ping under low load, and the corresponding parameters, like
throughout, \gls{IAT}, and packet size.
For some runs with \gls{OAI}, a disconnect during this stage was observed, which is also caught by the error handling routine.
Finally, temporary files are cleaned up and all produced artefacts are copied over.
The types of artefacts are explained in more detail in Section~\ref{sec:meth:artefacts}.





\subsection{Polling of Channel Metrics}\label{sec:meth:channelmetrics}

For the modem used in this work, the SIM8200EA-M2, it is possible to access multiple channel metrics directly
from the modem itself.
For this process, the QMI interface is used.
The process of extracting the metrics can be seen in Listing~\ref{lst:meth:qmiinfo}.

\gls{SNR}, \gls{RSRP}, and \gls{RSRQ} can be querried with the \texttt{-{}-nas-get-signal-info} flag and
\gls{SiNR} can be querried with the \texttt{-{}-nas-get-signal-strength} flag.
Only a limited amount of client IDs ca be created by QMI, therefore, the \texttt{ -{}-client-no-release-cid } option should be used.
After the first invocation, a client ID will be returned, which can be reused with \texttt{ -{}-client-cid=<ID> }.

In this setup, new values can be queried approximately every \SI{0.75}{\milli\second}.
To record SNR,RSRP, RSRQ, and SINR, the queries \texttt{-{}-nas-get-signal-info} and
\newline\texttt{-{}-nas-get-signal-strength} are alternated and the results are written to a CSV file.
Thereby resulting in new data points approximately every \SI{150}{\milli\second} per metric.

The use of \texttt{qmicli} causes the ModemManager to loose access to the modem, this is described in more detail in Section~\ref{sec:meth:issueqmi}.
Additionally, it is important to note, that this method of acquiring channel metrics is unstable.
Frequent timeouts and missing values are expected.
It is further expected for a significant portion of the runs to not provide any channel metrics from the modem.


\begin{lstfloat}[htb]
\begin{lstlisting}[basicstyle={\ttfamily\small\color{black}},style=base,caption={Extracting channel metrics from a modem.},label={lst:meth:qmiinfo}]
@$@ qmicli -d /dev/cdc-wdm0 --nas-get-signal-info --client-no-release-cid \
        --client-cid 5
[/dev/cdc-wdm0] Successfully got signal info
5G:
        RSRP: '-76 dBm'
        SNR: '29.5 dB'
        RSRQ: '-11 dB'
[/dev/cdc-wdm0] Client ID not released:
        Service: 'nas'
            CID: '5'

@$@ qmicli -d /dev/cdc-wdm0 --nas-get-signal-strength --client-no-release-cid \
        --client-cid 5
[/dev/cdc-wdm0] Successfully got signal strength
Current:
        Network '5gnr': '-65 dBm'
IO: '-106 dBm'
SINR (8): '9.0 dB'
[/dev/cdc-wdm0] Client ID not released:
        Service: 'nas'
            CID: '5'
\end{lstlisting}
\end{lstfloat}



\subsection{Measurement Artefacts}\label{sec:meth:artefacts}

For each run, whether successfull or not, a number of artefacts is collected.
Most importantly, each successfull test yields two PCAP recordings.
One containing all traffic captured at the modem and the other one containing all traffic captured at the Ethernet interface.
Additionally, the log file from the \gls{gNB} is saved which can be parsed to obtain channel metrics from the
point of view of the \gls{gNB}.
Tests employing the srsRAN \gls{gNB} additionally produce MAC layer PCAPs, which can be used to investigate
retransmissions or problems in the establishment of connections.
For \gls{OAI}, this functionality is not available.
There are also artefacts only produced by failed runs.
If a failure can be detected during execution of the playbook, the fault reason is saved.
Reasons for a fault for example include failure of the \gls{gNB} to start or the inability to acquire an IP address from the \gls{5GC}.
Further, for all runs, the output of all executed scripts and programms is saved.
These include the output of \texttt{iperf3}, the script used to establish a connection, the activation of performance tuning,
or the output of \texttt{uhubctl}.
Finally, for each test, the test definition which includes all parameters and configuration options for this run,
is also stored.
The extend of artefacts allows the operators to identify and explain outliers or faulty configurations.



\subsection{Explanation of Parameters}\label{sec:meth:parameters}

In this section, the utilized parameters are presented.
A general overview over the parameters can be seen in Table~\ref{table:meth:parameters}.
Appart from two specific scenarios, the distance between \gls{UE} and \gls{RU} is fixed at \SI{1.5}{\meter}
and a distance to nearby walls and objects is kept at \SI{1}{\meter}.
With the exception of the distance and the clearance, all parameters listed can be defined in the test definition
changed automatically via the Ansible playbook.

The \gls{UHD} is relevant for the communication between the \gls{gNB} and the \gls{SDR}.
Additionally, the \gls{UHD} is required for the compilation of the \gls{gNB}.
Therefore, in the setup phase of the testbed, each combination of \gls{UHD} version and \gls{gNB} version is
compiled once.
During measurement runs, it is only necessary to install the wanted \gls{UHD} version, i.e. copy the compiled libraries.
Then the correct \gls{gNB} binary for the required \gls{gNB} version which has been compiled against the current
\gls{UHD} version can be executed.

As explained in Section~\ref{sec:back:frame} and Section~\ref{sec:back:duplexing}, \gls{TDD} is utilized as the duplexing mechanism.
For this work, a \gls{TDD} pattern is defined by the period length and the ratio of downlink to uplink slots.
A listing of all utilized patterns can be seen in Table~\ref{table:meth:tddpattern}.
For two of those patterns, no \gls{OAI} measurements can be performed, since \gls{OAI} either crashes or is
unable to provide an IP address to the modem.

Performance tuning entails a combination of system tuning recommendations for srsRAN and \gls{OAI}.
In detail, when performance tuning is active, the kernel parameter \texttt{mitigations=off} is set,
the CPU governor is set to \texttt{performance}, DRM KMS polling is disabled, CPU idle C-states are disabled and
the network buffer size is increased.
The kernel parameter disables mitigations for CPU vulnerabilities, such as Spectre.
Changing the governor to \texttt{performance}, as well as disabling certain C-states forces the CPU to run at the maximum frequency,
and disabling DRM KMS polling prevents an operation in the Linux direct rendering manager, which sometimes helps performance.

There are two different types of traffic generation used throughout this work, throughout measurements with iperf and
ping measurements with a custom python script, utilizing scapy.
The traffic generation is described in depth in Section~\ref{sec:meth:trafficgen}.

Dockerization controls whether the precompiled \gls{gNB} binaries, as explained above, should be executed, or
a docker container for selected combination of \gls{UHD} version and \gls{gNB} version should be used instead.
Docker container for all combinations of \gls{UHD} version and \gls{gNB} version were built during the setup
phase of the testbed.



\begin{table}[htb]
\centering
\caption{Parameters and default values.}
\label{table:meth:parameters}
\begin{tabular}{l | l}
    \toprule
    Parameter & Manifestations \\
    \midrule
    Distance between \gls{UE} and core & \texttt{1.0m, \underline{1.5m} 2.0m} \\
    Clearance to wall & \texttt{\underline{1.0m}, 0.15m} \\
    \gls{UHD} Version & \texttt{3.15, \underline{4.0}} \\
    \gls{gNB} type & \texttt{srsRAN, OAI} \\
    \gls{gNB} version & srsRAN: \texttt{24.04, \underline{24.10}}, OAI: \texttt{2.1.0, \underline{2.2.0}} \\
    \gls{TDD} DL:UL ratio & \texttt{1, 2, 4} \\
    \gls{TDD} period & \texttt{5, 10, 20} \\
    Performance tuning & \texttt{\underline{true}, false} \\
    Traffic type & \texttt{iperf throughput, scapy ping} \\
    Traffic pattern & \texttt{\textit{scenario dependent}} \\
    Dockerization & \texttt{true, \underline{false}} \\

\end{tabular}
\end{table}

\begin{table}[htb]
\centering
\caption{TDD slot patterns.}
\label{table:meth:tddpattern}
\begin{tabular}{l | l | l | l | l | l | l}
    \toprule
    TDD period & $\frac{\text{DL}}{\text{UL}}$ ratio & Dl slots & Dl symbols & Ul slots & Ul symbols & $\mu$ \\
    \midrule
    20 slots      & 1:1 & 10 & 5 & 9 & 7 & 1 \\
    20 slots      & 2:1 & 13 & 5  & 6 & 7 & 1 \\
    20 slots$^{\color{Cblueish2}a}$  & 4:1 & 15 & 12 & 4 & 0 & 1 \\
    10 slots$^{\color{Cblueish2}ab}$ & 1:1 & 5  & 5  & 4 & 7 & 1 \\
    10 slots      & 2:1 & 6  & 8  & 3 & 4 & 1 \\
    10 slots      & 4:1 & 7  & 12 & 2 & 0 & 1 \\
    5 slots       & 1:1 & 2  & 6  & 2 & 6 & 1 \\
    5 slots       & 2:1 & 3  & 5  & 1 & 7 & 1 \\
    5 slots       & 4:1 & 3  & 12 & 1 & 0 & 1 \\
    \multicolumn{7}{l}{\footnotesize {\color{Cblueish2}${}^a$} OAI is not able to handle this configuration} \\
    \multicolumn{7}{l}{\footnotesize {\color{Cblueish2}${}^b$} Not a valid \gls{TDD} pattern~\cite[p. 90]{3gpp38.213},
    however the same patterns as in~\cite{raffeck2024parameterizing} are used.}
\end{tabular}
\end{table}

\subsection{Traffic Generation}\label{sec:meth:trafficgen}

Throughout this work, two forms of traffic generation are applied.
For throughput measurements, \texttt{iperf3} is used.
The \texttt{iperf3} client is executed on the \gls{gNB} host and \texttt{iperf3} in server mode is executed on
the \gls{5GC}. For downlink measurements, the \texttt{iperf3} client is executed with the reverse flag.
The generated bandwidth depends on the parameters of the run and is passed to the corresponding \texttt{iperf3} flag.
The protocol of all generated traffic is UDP.
The packets produced by \texttt{iperf3} contain sequence numbers, which are later used to match packets from
the modem interface and the Ethernet interface and calculate the latency of the system under load.

For delay measurements, a custom python script is used.
This script utilizes scapy to generate packets.
The script is used to generate a low volume of packets, significantly below the offered throughput of the system.
These packets are used to calculate the delay of packets when the system is under minimal load.
UDP packets with a sequence number are sent to a python server, running at the host of the \gls{5GC}.
The server responds with a packet containing the same sequence number, similar to an ICMP ping.
For this type of traffic generator, \gls{IAT} and packet size are used as parameters.


\subsection{Powercycle USB modem}\label{sec:meth:modemcycle}
To recreate a reproducible state and prevent remaining fragments of an old configuration, the modem is power cycled before each measurement run.
For this procedure an USB hub compatible with \texttt{uhubctl}\cite{uhubctl} is used.
\texttt{uhubctl} is a software which allows the selective up and down powering of specific ports on an USB hub.
The modem used for the measurements is the SIM8200EA-M2 5G HAT, which has seperate data and power USB connectors.
Both are connected to the USB Hub.
To prevent any saved state, \texttt{usb\_resetter}\cite{usbresetter} is used to rebind the complete USB hub
after the ports are powered down.
The combined use of \texttt{uhubctl} and \texttt{usb\_resetter} prevents a problem which is explained in more
detail in Section~\ref{sec:meth:probmodemcyclecache}.
In some cases, the USB hub is not detected by \texttt{uhubctl}, this issue is coverd in
Section~\ref{sec:meth:probmodemcycle}.



\section{Known Issues and Their Mitigation}\label{sec:meth:issues}

During the development of the automation framework, several behaviors and system faults were encountered.
These technical issues affect the system stability or measurement reliability.
This section describes the most significant problems observed during testing.
Where possible, mitigation procedures are provided.

\subsection{\gls{SDR} Fallback to USB2 Mode}

After many test runs (approximately 1 in hundreds), it is possible for the B210 \gls{SDR} to fall back
from USB3 to USB2 connectivity.
When the \gls{SDR} is in USB2 mode, srsRAN will exit with failure after a few seconds and OAI
will start but not provide any IP addresses.
Rebinding the \gls{SDR} in software was not successfull.
To set the \gls{SDR} back to USB3 mode, it is required to detatch the usb cable and powercycle the \gls{SDR}.

\texttt{lsusb -t} can be used to show the speed.
The \gls{SDR} should be running at \SI{5000}{\mega\byte\per\second}.
When the \gls{SDR} is powercycled, it will only show present one USB device instead of 4 and list a speed of 480 which indicates USB2 mode.
After the firware is flashed to the \gls{SDR} by the \gls{UHD} driver, it will present as 4 USB devices with the updated speed.

It was not possible to determine the source of this behavior, however, during all test runs, this fault only appeared
after a measurement with OAI.
The mitigation of this behavior was not automated in this work.
However, in the future, a smart plug could be used to integrate power cycling the \gls{SDR} into the Ansible playbook for measurements.


\subsection{Caveats of Using qmicli}\label{sec:meth:issueqmi}

When using the methodology described in \ref{sec:meth:channelmetrics}, it is important to
perform all required \texttt{ModemManager} operations beforehand.
After the use of \texttt{qmicli}, \texttt{ModemManager} will loose access to the device.
Additionally, it has to be kept in mind, that that this method of acquiring channel metrics from the modem is not
reliable and frequent timeouts and missing values are expected.

\subsection{Inconsistent Interface Name After a Reboot}\label{sec:meth:issueinterfacename}

After some reboots, the interface name of the USB modem changes unexpectedly.
The nameing scheme changes in this case from \texttt{wwan0} to a MAC based naming scheme, e.g. \texttt{wxn01ab23cd45ef}.
This behavior can be observed in approximately \SI{20}{\percent} of all cases.
To resolve this issue, another reboot can be performed.

\subsection{Failure to Power Cycle the Modem}\label{sec:meth:probmodemcycle}

Periodically, after a run with OAI, the connection to the USB modem becomes broken and a power cycle
via the method described in Section~\ref{sec:meth:modemcycle} is no longer possible.
To prevent undefined behavior, a reboot of the host machine is performed.
This problem is detected and handled automatically as explained in Section~\ref{sec:meth:measureautomation}.

\subsection{Failure to Invalidate Cache After Modem Power Cycle}\label{sec:meth:probmodemcyclecache}

In some cases, after the connected ports on the USB hub are turned off by \texttt{uhubctl}, the modem is stil
listed as a connected device in the operating system.
Additionally, \emph{(cached)} information about the SIM can be read by the ModemManager as if it were still powered on.
Physically disconecting the modem does not change this behavior.
The saved state could be a source of unintentional influences.
Therefore, \texttt{usb\_resetter} is used in conjunction with \texttt{uhubctl} to temporarly disable
and rebind the USB hub after the relevant ports of the hub have been powered down by \texttt{uhubctl}.
During the unbind, the USB hub and all connected devices are no longer accessible and the cached reference to
the modem is cleaned up.

\subsection{Unexpected Halts in OAI Measurement Runs}

In some rare cases, an automated measurement can freeze completely.
This only happend three times during all measurements conducted for this work.
The problem also only occured when performing measurements with OAI as \gls{gNB}.
The source of this problem could not be isolated.
Therefore no steps to mitigate this issue were taken.




\chapter{Evaluation}\label{chap:evaluation}

The evaluation of this work is split into two parts.
First, the influences of individual parameters are examined in a collection of scenarios.
This step isolates the effects of specific parameters and helps build a foundation for interpreting more complex interactions.
The foundation is used to design a comprehensive case study, which applies a full factorial experiment across a broader set of parameters.
Through a full factorial experiment, the case study aims to identify which parameters and parameter combinations
are most critical to achieving reliable and reproducible results in 5G campus deployments.



\section{Setup validation}\label{sec:eval:validation}
In this section, the measurement scenarios are presented.

First, the configuration of the transmission and reception gain is expected to have a significant impact on the
quality of the transmission channel.
This factor is not mentioned in the covered related work, yet typically is critical parameter for improving the signal strength.
Therefore, the influence of the gain configuration is investigated in Section~\ref{sec:eval:gain}.

In \cite{raffeck2024parameterizing}, a significant impact of the elevation between the \gls{UE} and \gls{RU} was reported.
Thus, influences of the positioning of \gls{UE} and \gls{RU} are analyzed in Section~\ref{sec:eval:positioning}.
The positioning includes the elevation between \gls{UE} and \gls{RU}, the distance between them, as well as
the clearance to nearby walls.

For srsRAN and OAI, system settings for optimized performance are recommended.
However in both cases, the recommendations are given with the caveat, that the actual performance must be examined on a
case-by-case basis.
The configuration recommendations for OAI and srsRAN significantly overlap, therefore, in
Section~\ref{sec:eval:perftuning}, the impact of these configuration recommendations are investigated.

One important trait of a 5G campus deployments is the software version of the \gls{gNB} application and the driver.
Since each application version is compiled against the current driver version, which in turn is compiled against
the Linux header files, compilation is time consuming and produces binaries which can not easily be
swapped between different systems.
To ease the testing of different application and driver versions, the feasibility of dockerization of these
applications is investigated in Section~\ref{sec:eval:docker}.

Throughput measurements are conducted by generating a data volume higher than the system is able to transmit.
In the scenario Section~\ref{sec:eval:throughput}, throughput measurements are performed with a differing amounts
of generated data volume.
This scenario enables two key observations.
First, it reveals whether the system exhibits different behavior depending on the degree to which the offered traffic
exceeds the maximum throughput, for instance comparing \SI{110}{\percent} to \SI{200}{\percent} of the system capacity.
Second, it provides insight how the system behaves when operating under high, yet non-saturating, load compared to full capacity.

The modulation and code rate is selected based on the channel condition and a \gls{MCS} table.
Depending on the maximum \gls{QAM} level, different tables are available.
While all other measurements are performed with a maximum modulation of 256\gls{QAM}, srsRAN allows to
set 64\gls{QAM} as maximum modulation level.
The effect of this difference in modulation is investigated in Section~\ref{sec:eval:qam}.

A broad selection of channel quality indicators and signal strength metrics are presented in Section~\ref{sec:back:channel}.
In Section~\ref{sec:eval:timeseries}, the metrics as reported by the \gls{gNB} are juxtaposed to the per-packet latency
during one sample measurement run.
This scenario aims to illustrate the value of the channel metrics and how they can provide better insights into random effects and measurement anomalies.


\subsection{Reception and Transmission Gain}\label{sec:eval:gain}

Before a valid measurement can be conducted, it is important to configure the transmission gain of the \gls{gNB}.
The method for this is described in Section~\ref{sec:meth:antennagain} and must be repeated after each
change in location or distance between \gls{gNB} and \gls{UE}.
In this section, the impact of different gain settings is investigated.

In Figure~\ref{fig:eval:gain.iperf.throughput} the offered throughput in relation to the type of \gls{gNB},
OAI or srsRAN, and the measurement direction, downlink or uplink, can be seen.
For downlink measurements, up to a certain point, an increased Rx gain, as well as an increased Tx gain, results
in an improved throughput.
The measured throughout rises up to \SI{50}{\mega\bit\per\second} for \gls{OAI} and \SI{41}{\mega\bit\per\second} for srsRAN.
If the maximum achievable throughput is reached, further increasing the gain does not result in inproved throughput.
For \gls{OAI}, the throughout plateaus for a reception gain of \SI{38}{\deci\bel} and a transmissions gain of \SI{76}{\deci\bel}.
The throughout for srsRAN downlink measurement plateaus at lower values, at a reception gain
of \SI{8}{\deci\bel} and a transmission gain of \SI{74}{\deci\bel}.
For uplink measurements, the same trend can be seen for OAI as \gls{gNB}.
Showing an increase in throughput for a higher gain, up to approximately \SI{18}{\mega\bit\per\second}.
However, in contrast to downlink measurements, no plateau is reached.
Uplink measurements with srsRAN, on the other hand, display a different behavior.
The measured throughput also increases up to certain point.
Further increasing the gain, however, leads to a decreased throughput performance.
This effect is particularly well visible for the uplink srsRAN measurement with varying Rx gain.
The peak uplink performance is at \SI{44}{\deci\bel} and \SI{74}{\deci\bel} for reception and transmission gain, respectively.
The maximum uplink throughput is at \SI{16}{\mega\bit\per\second}.

The positive impact of the gain on the throughput is expected.
However, to better understand the cause of the adverse effect of high gain values, the \gls{SNR} as reported by the \gls{gNB} is plotted in
Figure~\ref{fig:eval:gain.iperf.snrgnb}.
The same parameter combinations as before are used, facetted for the type of \gls{gNB} and the traffic direction.
The reported \gls{SNR} values from OAI show very little variation and are with a high likelyhood invalid.
In contrast, the \gls{SNR} values reported by srsRAN show similar behavior for uplink and downlink, as well as for reception and transmission.
There is a local maximum, significantly increasing or decreasing the gain setting results in an decrease of \gls{SNR}.

The \gls{SNR} as reported from the modem instead of the \gls{gNB} can be seen in Figure~\ref{fig:eval:gain.iperf.snrmodem}.
Due to the limitations explained in Section~\ref{sec:meth:issueqmi}, many datapoints are missing,
especially for measurements conducted with \gls{OAI} as \gls{gNB}.
In all cases, the reception gain seems independent from the \gls{SNR}, displaying about \SI{30}{\deci\bel} for srsRAN and \SI{27}{\deci\bel} for OAI.
This is expected behavior, since the amplification at the reception antenna of the \gls{gNB} is not visible to the modem.
The \gls{SNR} for the transmission gain configuration, on the other hand, rises with an increase in gain.
This effect is well visible for srsRAN.
A similar pattern appears to emerge for \gls{OAI}, but due to the incomplete dataset, this conclusion should be treated with caution.

The findings from the reported \gls{SNR} further substantiate the most likely explanation for this behavior.
An increased gain amplifies the signal, therefore it can be more easily distinguished from background noise.
Simultaneously, a signal causes reflections, which grow unproportional higher, the more amplified the signal is.
Additionally, a higher degree of gain also amplifies noise and thereby creates more errors during the digitization of an analog signal at the reception antenna.
There is a sweet spot for those effects at which the optimal system performance under the given circumstances can be achieved.
For the current environment and configuration, this sweet spot is at a Rx gain of \SI{44}{\deci\bel} and a Tx gain of \SI{74}{\deci\bel} for srsRAN.
In the case of \gls{OAI}, the best performance can be achieved at a Rx gain of \SI{68}{\deci\bel} and a Tx gain of \SI{88}{\deci\bel}.


The inverse effect can be seen for the \gls{OWD} in Figure~\ref{fig:eval:gain.iperf.delay}.
The observed delay behaves inversely proportional to the throughput.
A lower delay directly increases the time a packet spends in the network interface buffer,
significantly increasing the observed delay.


\begin{figure}[htb]
\centering
\includegraphics[width=0.95\textwidth]{/mnt/ext1/5g-masterarbeit-daten/antenna-gain-b205/rxtx-gain-throughput-agg_alt01.pdf}
\caption{Difference in throughput for a range of Rx and Tx antenna gain values.}
\label{fig:eval:gain.iperf.throughput}
\end{figure}

\begin{figure}[htb]
\centering
\includegraphics[width=0.95\textwidth]{/mnt/ext1/5g-masterarbeit-daten/antenna-gain-b205/rxtx-gain-gnb_snr-agg_alt01.pdf}
\caption{\gls{SNR} at the \gls{gNB} for varying Rx and Tx gain settings.}
\label{fig:eval:gain.iperf.snrgnb}
\end{figure}

\begin{figure}[htb]
\centering
\includegraphics[width=0.95\textwidth]{/mnt/ext1/5g-masterarbeit-daten/antenna-gain-b205/rxtx-gain-modem_snr-agg_alt01.pdf}
\caption{\gls{SNR} at the modem for varying Rx and Tx gain settings.}
\label{fig:eval:gain.iperf.snrmodem}
\end{figure}


\begin{figure}[htb]
\centering
\includegraphics[width=0.95\textwidth]{/mnt/ext1/5g-masterarbeit-daten/antenna-gain-b205/rxtx-gain-delay-agg_alt01_zoom.pdf}
\caption{Difference in delay for a range of Rx and Tx antenna gain values.}
\label{fig:eval:gain.iperf.delay}
\end{figure}



\subsection{Distance and Height}\label{sec:eval:positioning}

In \cite{raffeck2024parameterizing}, a significant decrease in throughput for a different elevation of \gls{gNB} and
\gls{UE} can be seen.
Therefore, in this section, the influence of the positioning of \gls{gNB} and \gls{UE} is investigated.
The calibration measurement for this scenario has been performed once for a distance of \SI{1.5}{\meter}.

In Figure~\ref{fig:eval:distance.iperf.throughput}, the offered throughput, in relation to the \gls{TDD} period length,
for horizontal facets, and the \gls{gNB} type and traffic direction for vertical facets can be seen.
The \gls{TDD} ratio is color coded and the height is noted on the x-axis, while the y-axis shows the throughput in \SI{}{\mega\bit\per\second}.
For almost all cases, no measured effect can be seen.
For OAI uplink measurements, as well as srsRAN uplink measurements for a \gls{TDD} pattern of 5 slots and a ratio
of 1:1 and a \gls{TDD} pattern of 10 slots and a ratio of 4:1, a decreased mean throughput performance can be seen
for \SI{1.5}{\meter}, while the results for \SI{1}{\meter} and \SI{2}{\meter} are near identical.
This effect, however, is not statistically significant.

These results are unanticipated and a decrease in throughput for a higher distance would have been expected.
Therefore, the reported \gls{SNR} of the modem for the same parameter combinations is investigated in Figure~\ref{fig:eval:distance.iperf.modemsnr}.
As explained in Section~\ref{sec:meth:issueqmi}, due to the unreliability of the metric polling methodology, multiple data points are missing.
Nonetheless, a clear trend can be seen for srsRAN downlink and uplink measurements.
All aquired values show a significant better \gls{SNR} for a smaller distance than for a greater distance,
decreasing from just over \SI{29}{\deci\bel} for a distance of \SI{1}{\meter} to approximately \SI{27.3}{\deci\bel}
at a distance of \SI{1.5}{\meter} and \SI{26}{\deci\bel} at a distance of \SI{2}{\meter}.
For OAI measurements, not enough values are a available to infer an explanation.
The reported values for the srsRAN measurements, however, indicate, that for this scenario, the bottleneck of the throughput is not the channel quality.
Despite a decline in channel quality with an increasing distance, no corresponding worsing of throughput can be observed.

\begin{figure}[htb]
\centering
\includegraphics[width=0.95\textwidth]{/mnt/ext1/5g-masterarbeit-daten/distance/agg_performance_tuning_iperf-throughput-c.pdf}
\caption{Influence of distance between \gls{RU} and \gls{UE} on the offered throughput, dependent on \gls{TDD} pattern and \gls{gNB} type.}
\label{fig:eval:distance.iperf.throughput}
\end{figure}

\begin{figure}[htb]
\centering
\includegraphics[width=0.95\textwidth]{/mnt/ext1/5g-masterarbeit-daten/distance/agg_performance_tuning_iperf-modemsnr-c.pdf}
\caption{Influence on distance on the \gls{SNR} as reported by the modem.\todo{Plot stauchen oder aggregieren}}
\label{fig:eval:distance.iperf.modemsnr}
\end{figure}


\begin{figure}[htb]
\centering
\captionsetup{width=0.46\textwidth}
\begin{minipage}[t]{0.49\textwidth}
    \centering
    \includegraphics[width=0.99\textwidth]{/mnt/ext1/5g-masterarbeit-daten/distance_wall/agg_performance_tuning_iperf-throughput.pdf}
    \caption{Influence on clearance to nearest wall on offered throughput for an iperf measurement.}
    \label{fig:eval:clearance.iperf.throughput}
\end{minipage}%
\begin{minipage}[t]{0.49\textwidth}
    \centering
    \includegraphics[width=0.99\textwidth]{/mnt/ext1/5g-masterarbeit-daten/distance_wall/agg_performance_tuning_scapy-delay.pdf}
    \caption{Influence on clearance to nearest wall on delay for a ping measurement.}
    \label{fig:eval:clearance.scapy.delay}
\end{minipage}
\end{figure}

Throughout this work, the \gls{RU} and the \gls{UE} always maintain a clearance of at least \SI{1}{\meter} to nearby
walls and other reflective surfaces.
The effect of ignoring this minimal clearance can be seen in Figure~\ref{fig:eval:clearance.iperf.throughput}
and Figure~\ref{fig:eval:clearance.scapy.delay}.
For this scenario, the tripods of the \gls{RU} and \gls{UE} are placed directly next to a nearby concrete wall
resulting in a clearance of \SI{15}{\centi\meter} between wall and antenna for both devices.
In Figure~\ref{fig:eval:clearance.iperf.throughput} the effect on the throughput and in Figure~\ref{fig:eval:clearance.scapy.delay}
the effect on the delay can be seen.
In both cases, only a \gls{TDD} configuration with a length of five slots and an downlink to uplink ratio of 1:1 was inspected.
Different \gls{gNB} types and traffic directions are differentiated by the vertical facets.
In regards to the throughput measurement, only on statistically significant observation can be made.
The uplink throughout of OAI is significantly better near the wall than with a higher clearance.
The same observation can be made for the srsRAN uplink measurement, however, the effect is not statistically significant.
In general, a trend towards increased throughout by reducing the clearance can be seen.
For the delay measurement of ping type traffic in Figure~\ref{fig:eval:clearance.scapy.delay}, no significant effect can be seen.
However, a small trend towards a slightly increased delay for a reduced clearance is visible.



\begin{figure}[htb]
\centering
\includegraphics[width=0.95\textwidth]{/mnt/ext1/5g-masterarbeit-daten/height/agg_performance_height_iperf-throughput.pdf}
\caption{Influence of \SI{1}{\meter} height difference between \gls{gNB} and \gls{UE}}
\label{fig:eval:height.iperf.throughput}
\end{figure}
\begin{figure}[htb]
\centering
\includegraphics[width=0.95\textwidth]{/mnt/ext1/5g-masterarbeit-daten/height/agg_performance_height_iperf-modemsnr.pdf}
\caption{Influence of \SI{1}{\meter} height difference on the \gls{SNR} as reported by the modem.\todo{Plot stauchen oder aggregieren}}
\label{fig:eval:height.iperf.modemsnr}
\end{figure}

In \cite{raffeck2024parameterizing} a significant reduction in offered throughput was reported for a configuration with an elevated \gls{gNB}.
Therefore, the influence of height differences between \gls{RU} and \gls{UE} is investigated in Figure~\ref{fig:eval:height.iperf.throughput}.
For this measurement scenario, the distance between \gls{RU} and \gls{UE} is kept at the default \SI{1.5}{\meter}.
The offered throughput in \SI{}{\mega\bit\per\second} is indicated on the y-axis and the hight difference is marked on the x-axis.
A positive value corresponds to an elevated \gls{RU}.
The values are facetted for \gls{TDD} period length, traffic direction, and \gls{gNB} type.
The \gls{TDD} ratio is indicated by the color.
For one parameter combination, a period length of five slots and uplink traffic with the OAI \gls{gNB},
a minor but statistically significant reduction in offered throughput can be seen.
In all other cases, no effect can be observed.
Therefore, likewise to the distance measurements, the \gls{SNR} of the modem was investigated.
The results can be seen in Figure~\ref{fig:eval:height.iperf.modemsnr}, with the same facets and color coding as the figure for the throughput measurement.
Similarly to the behavior of the \gls{SNR} in Figure~\ref{fig:eval:distance.iperf.modemsnr},
while clear no effect can be seen in regards to the throughput measurement,
the influence of the parameter on the \gls{SNR} can be reliably observed, despite many data point missing.
For srsRAN measurements, the measured \gls{SNR} decreases from approximately \SI{30}{\deci\bel} to \SI{22.5}{\deci\bel} when \gls{RU} and \gls{UE} are at different heights.
In the same way, the \gls{SNR} for OAI measurements decreases from approximately \SI{27.5}{\deci\bel} to \SI{24.5}{\deci\bel}.
Therefore, while the immediate effect of \cite{raffeck2024parameterizing} could not be reproduced, a strong adverse effect of the
height difference between \gls{RU} and \gls{UE} can be observed.
This observation further amplifies the importance of correct adjustment of antenna gain,
since a bad or missing calibration could easily result in a system directly bottlenecked by the channel quality.
If the system is bottlenecked by the channel quality, adverse influences, like an increased distance or height difference
between \gls{RU} and \gls{UE}, would directly affect the observed maximum throughput.



\subsection{Performance Tuning}\label{sec:eval:perftuning}

For both types of \gls{gNB} application used in this work, srsRAN and OAI, host configuration recommendations for a
better performance are given.
In this section, the effect of the given recommendations are evaluated.
Perfromance recommendations are combined to form a single parameter, performance tuning or no performance tuning.
The configuration steps this parameter is comprised of are explained in Section~\ref{sec:meth:parameters}.
The effect of this parameter on the system throughput is highlighted in Figure~\ref{fig:eval:performancetuning.iperf.bandwidth}.
There are horizontal and vertical facets for the traffic direction and the \gls{gNB} type respectively.
The \gls{TDD} configuration is marked at the x-axis and the offered throughput is marked at the y-axis.
There are no significant effects for uplink traffic. For downlink traffic, there are no significant effects for
srsRAN, however OAI downlink a small but significant improvement for all tested configurations.
In Figure~\ref{fig:eval:performancetuning.iperf.delay}, the delay of the same parameter configurations can be seen.
The uplink delay for both \gls{gNB} types is significantly increased for the tuning group.
This effect can be explained by the increased buffer sizes, resulting in a longer queuing delay at the same throughput.
The \gls{OWD} for downlink transmission with OAI is reduced by performance tuning, likely as a result of
the increased throughput.
Interestingly, \gls{OWD} for downlink transmissions with srsRAN is not affected in any significant way.
The \gls{OWD} for ping measurements, under no system load, show no significant effect of performance tuning.
This observation affirms, that performance tuning can have an impact, however, a more fine-grained analysis is required
to reliably improve system performance while preventing unintended regressions.
The splitting of this observation into multiple parameters, like changing the governor, setting the buffer sizes,
and preventing certain system calls, is reserved for future work.

\begin{figure}[htb]
\centering
\includegraphics[width=0.95\textwidth]{/mnt/ext1/5g-masterarbeit-daten/performance-tuning_cstate_recommends/agg_performance_tuning_iperf-bandwidth_alt01.pdf}
\caption{Difference in throughput between no tuning and performance tuning for iperf throughput measurements.}
\label{fig:eval:performancetuning.iperf.bandwidth}
\end{figure}

\begin{figure}[htb]
\centering
\includegraphics[width=0.95\textwidth]{/mnt/ext1/5g-masterarbeit-daten/performance-tuning_cstate_recommends/agg_performance_tuning_iperf-delay_alt01.pdf}
\caption{Difference in delay between no tuning and performance tuning for iperf throughput measurements.}
\label{fig:eval:performancetuning.iperf.delay}
\end{figure}



\subsection{Dockerization}\label{sec:eval:docker}

To change the application version of OAI or srsRAN, the application must be compiled against the selected
\gls{UHD} version, which in turn must be compiled against the Linux kernel headers.
This process is time consuming and hinders the evaluation of performance differences caused by different
\gls{UHD} or \gls{gNB} versions.
Therefore, for this section, we built Docker container for the product of the \gls{UHD} and \gls{gNB} versions used
throughout this work.

In Figure~\ref{fig:eval:dockerization.iperf.throughput} the result of a throughput measurement can be seen.
Horizontal facets differentiate between the OAI and srsRAN versions and vertical facets contrast between
the transmission direction, as well as the \gls{TDD} downlink to uplink ratio.
The \gls{TDD} period is color coded. On the y-axis the offered throughput in \SI{}{\mega\bit\per\second} is shown
and the x-axis displays whether the deployed \gls{gNB} was running on bare metal or in a docker container.
Both versions of srsRAN show no significant effect in all tested configurations and show the same system performance
for dockerized and not dockerized deployments.
On the other hand, both versions of OAI display a strong drop in throughput performance, loosing up to \SI{60}{\percent}
of downlink throughput and almost all uplink throughput.
As a consequence, dockerized OAI deployments show drastic increases in delay under load, up to \SI{20}{\second} for
downlink transmissions and up to \SI{30}{\second} for uplink transmissions.
The delay for ping measurements, with small packets, high \gls{IAT} and therefore low system load, can be seen in
Figure~\ref{fig:eval:dockerization.scapy.ping}.
While no significant effect between dockerization and no dockerization can be observed for all
evaluated srsRAN configurations, the dockerized OAI deployments display significant increased delay.
There are 3 downlink configurations for OAI which show no significant effect, for OAI 2.1.0 and 2.2.0 with a
\gls{TDD} period of 20 slots and a ratio of 2:1, and only for OAI 2.2.0 with a \gls{TDD} period of 10 slots and a ratio of 2:1.
All other OAI downlink configurations display significant increase in \gls{OWD} in the range of \SI{0.5}{\milli\second}
to \SI{2}{\milli\second}.
Many uplink transmission configurations of OAI feature an increased delay in the range of multiple seconds, since
the offered throughput is not enough to reliable transmit the periodic ping-type packets at a low system load.
This results in filling network buffers and a delay of multiple seconds, similar as in the throughput case.

The expected outcome of this scenario would have been a small but significant overhead, resulting in a slight
decrease in offered throughput for both types of \gls{gNB}.
Instead, no effect could be observed for srsRAN and a significant degradation of OAI performance could be seen.
OAI potentially depends on some type of hardware feature, which is not provided by docker, resulting in
the observed performance impact.


\begin{figure}[htb]
\centering
\includegraphics[width=0.95\textwidth]{/mnt/ext1/5g-masterarbeit-daten/dockerization/agg_performance_tuning_iperf-throughput-alt02.pdf}
\caption{Difference in throughput between dockerization and bare deployments for iperf throughput measurements.}
\label{fig:eval:dockerization.iperf.throughput}
\end{figure}

\begin{figure}[htb]
\centering
\includegraphics[width=0.95\textwidth]{/mnt/ext1/5g-masterarbeit-daten/dockerization/agg_performance_tuning_scapy-delay-alt02.pdf}
\caption{Difference in delay between dockerization and bare deployments for scapy ping measurements.}
\label{fig:eval:dockerization.scapy.ping}
\end{figure}



\subsection{Offered Throughput and Offered Load}\label{sec:eval:throughput}

Throughput measurements in this work are conducted by generating a higher data volume, than the system is able to transmit.
In this scenario, the effect of the generated data volume on the actual achieved throughput is investigated.
In Fiure~\ref{fig:eval:throughputovershoot.iperf.bandwidth}, the results of this investigation are visualized.
The magnitude of the generated traffic is varied along the x-axis and the system throughput is marked on the y-axis.
Horzontal facets and vertical facets differentiate the \gls{TDD} ratio and the type of \gls{gNB}, respectively.
The \gls{TDD} period is color coded and the direction is indicated by a solid or a dashed line.

For all srsRAN configurations, expected behavior can be observed.
As long as the system is able to handle the data volume, the throughput increases linearly.
At some point the system is no longer able the generated data volume and the throughput plateaus.
For multiple OAI configurations, most pronounced for the downlink configuration with a \gls{TDD} ratio of 2:1
and a \gls{TDD} period of 5, a different behavior can be observed.
At maximum load, the system throughput is significantly less, than at a lower load.
A possible explanation of this behavior is that the release of reserved ressources, i.e. if a
transmission must be repeated but the queue is full, leads to expensive operations which in turn negatively
affect the throughput.
Additionally, no statistically significant influence can be observed across all configurations with respect to how much the generated traffic exceeds the maximum throughput.
Noteably, the parameter defining the offered traffic rate does not acurately represent the actual load on the system.
Once the transmission buffer of the interface is saturated, the associated socket no longer accepts additional data and the traffic generation is throttled.
As a result, when the system is able to transmit up to \SI{70}{\mega\bit\per\second}, offered traffic configurations of \SI{100}{\mega\bit\per\second},
\SI{150}{\mega\bit\per\second}, or \SI{200}{\mega\bit\per\second} provide no significant difference to the data volume generated or transmitted.

\begin{figure}[htb]
\centering
\includegraphics[width=0.95\textwidth]{/mnt/ext1/5g-masterarbeit-daten/throughput-overshoot/agg_throughput_compare-bandwidth_169.pdf}
\caption{Influence of backpressure on the achieved throughput for iperf throughput measurements.}
\label{fig:eval:throughputovershoot.iperf.bandwidth}
\end{figure}


\subsection{QAM}\label{sec:eval:qam}

For srsRAN, it is possible to change the maximum \gls{QAM} level, and thereby the underlying \gls{MCS} scaling table,
in the configuration file.
In this scenario the difference between 64\gls{QAM} and 256\gls{QAM} for srsRAN release 24.10 is investigated.
In Figure~\ref{fig:eval:qam.iperf.bandwidth}, the effect of a more efficient modulation scheme can be seen.
The upper row displays results for downlink measurements and the lower row displays results for uplink measurements.
Horizontal facets and the color denote the \gls{TDD} period length and the \gls{TDD} ratio respectively.
The achieved throughput is marked at the y-axis and the \gls{QAM} level is indicated at the x-axis.

For all measurements, a clear and significant positive impact for a higher \gls{QAM} level is visible.
As explained in Section~\ref{sec:back:channelmcs}, a higher \gls{QAM} level typically results in an increased
\gls{BLER}, therefore, the best suitable modulation scheme, typically depends on the quality of the communication channel.
Simultaneously, a higher modulation level increases the use of spectral resources and thereby the transmissions capacity.
Under the circumstances of this measurement series, allowing a higher maximum \gls{QAM} level resulted in an increase
of \SI{23.5}{\percent} to \SI{42.1}{\percent} for uplink traffic and \SI{30.9}{\percent} to \SI{31.6}{\percent} for downlink traffic.

\begin{figure}[htb]
\centering
\includegraphics[width=0.95\textwidth]{/mnt/ext1/5g-masterarbeit-daten/qam/agg_performance_tuning_iperf-throughput-24_10-alt-169.pdf}
\caption{Throughput performance difference between 64\gls{QAM} and 256\gls{QAM}.}
\label{fig:eval:qam.iperf.bandwidth}
\end{figure}



\subsection{Time Series Analysis}\label{sec:eval:timeseries}

\begin{figure}[htb]
\centering
\includegraphics[width=0.95\textwidth]{/home/lks/DocSync/iat-channel--nomodem--zstd-20000_3_1_.pdf}
\caption{Time series of delay an channel metrics of an iperf3 measurement.}
\label{fig:eval:timeseries}
\end{figure}

In Section~\ref{sec:meth:channelmetrics}, the process of polling channel metrics is described.
In addition, to the collected data on the modem, \gls{gNB} itself reports several metrics, including \gls{CQI}, \gls{SNR}, \gls{RSRP}, and uplink and downlink \gls{MCS}.
The collected information of the channel quality can be consulted to explain measurement anomalies.
While in Section~\ref{sec:eval:positioning} the behavior of aggregated \gls{SNR} is discussed, this section illustrates how channel
quality metrics evolve during a single measurement run.

In Figure~\ref{fig:eval:timeseries}, a \SI{6}{\second} cut-out of a throughput measurement can be seen.
The per-packet delay in seconds is displayed in the upper facet and the channel metrics as reported by srsRAN are shown in the lower facet.
To better highlight deviations from the mean of the channel quality indicators, the values have been z standardized.
The resulting z-score is marked on the y-axis, with a score of zero representing the mean and a score of one corresponding to one standard deviation above the mean.
In the selected time window, a temporary increase in delay is observable.
During the onset and the initial phase of this increase, several outliers with an increased delay are visible.
These outliers are potentially caused by retransmissions.
As a result, the downlink \gls{MCS} is reduced, resulting in a lower throughput but also an increased robustness against transmission errors.
A reduction in throughput, on the other hand, also directly affects the queueing delay and thereby increasing the total observed delay.
Interestingly, during this period, an increase in the measured \gls{RSRP} is accompanied by a decrease in \gls{SNR}.
This is counterintuitive, since a higher \gls{RSRP} typically correlates with higher \gls{SNR}.
This unexpected behavior cannot be fully explained at this point.
When the \gls{MCS} is increased again, the delay first plateaus, after which it decreases to baseline.

This example highlights, how the collection of channel metrics and quality indicators can help in explaining
anomalies in the observed traffic.
Thereby, it serves an important role in providing a better explainability and reproducibility of the study.
However, some limitations have to be kept in mind.
First, as described in Section~\ref{sec:meth:issueqmi}, the modem's channel metrics could not be collected reliably.
Second, as shown in Section~\ref{sec:eval:gain}, OAI reports questionable \gls{SNR} values,
raising concerns about the trustworthiness of its channel metrics.
Finally, the relatively infrequent reporting interval of one second, used by both srsRAN and OAI,
limits the ability to observe fine-grained effects.



\section{Case study}\label{sec:eval:case}


In this section, a case study is presented to investigate the combined effects of key parameters on system performance.
This evaluation build upon the individual findings discussed previously.
To systematically analyze parameter interactions, a full factorial experimental design was selected.
However, there are some parameters, which are only available for specific traffic types.
The parameters and their manifestations are shown in Table~\ref{table:eval:parameters}.
In all cases, the \gls{UHD} version, the \gls{gNB} version and the \gls{TDD} pattern is included.
For throughput type measurements, the additional parameter direction is included.
Ping type measurements are already bidirectional, however, the additional parameters \gls{IAT} and packet size are included, which are not available for throughput measurements.

In this section, further effects are investigated.
First, in Section~\ref{sec:eval:cdf}, the \gls{CDF} of multiple runs is presented, and the approach of aggregating measurements, on a per-run basis is explained.
Then, the influence of different \gls{UHD} and \gls{gNB} versions are analyzed in Section~\ref{sec:eval:versions}.
The theoretical effect of the \gls{TDD} pattern is evaluated against the observed measurements in Section~\ref{sec:eval:tdd}.
In Section~\ref{sec:eval:cnsm}, the measurement results of the case study are compared to the reported measurements in \cite{raffeck2024parameterizing}.
Finally, for both types of measurements, a full factor \gls{ANOVA} is performed in Seciton~\ref{sec:eval:anova}.
The results are discussed, and the most prominent two-factor interactions are visualized in interaction plots.



\begin{table}[htb]
\centering
\caption{Parameters for case study.}
\label{table:eval:parameters}
\begin{tabular}{l | l | l}
    \toprule
    Parameter & Throughput & Ping delay \\
    \midrule
    \gls{UHD} version       & \texttt{3.15, 4.0}                & \texttt{3.15, 4.0} \\
    \gls{gNB} type$^{\color{Cblueish2}a}$     & \texttt{srsRAN, OAI}              & \texttt{srsRAN, OAI}  \\
    \gls{gNB} version       & srsRAN: \texttt{24.04, 24.10},    & srsRAN: \texttt{24.04, 24.10}, \\
                            & OAI: \texttt{2.1.0, 2.2.0}        & OAI: \texttt{2.1.0, 2.2.0} \\
    \gls{TDD} DL:UL         & \texttt{1, 2, 4}                  & \texttt{1, 2, 4} \\
    \gls{TDD} period        & \texttt{5, 10, 20}                & \texttt{5, 10, 20} \\
    Traffic type            & \texttt{iperf throughput}         & \texttt{scapy ping} \\
    Direction               & \texttt{DL, UL}                   & \textit{fixed} \\
    \gls{IAT}               & \textit{fixed}                    & \texttt{0.01, 0.001} \\
    Packet size             & \textit{fixed}                    & \texttt{36, 1382} \\
    \multicolumn{3}{l}{\footnotesize {\color{Cblueish2}${}^a$} Redundant for factorial design but included for evaluation} \\

\end{tabular}
\end{table}




\subsection{CDF of Delay}\label{sec:eval:cdf}

\begin{figure}[htb]
\centering
\captionsetup{width=0.46\textwidth}
\begin{minipage}[t]{0.49\textwidth}
    \centering
    \includegraphics[width=0.99\textwidth]{/mnt/ext1/5g-masterarbeit-daten/main_measurement/.plots/CDF_delay_9ff091db__580213e7_full_4:3.pdf}
    \caption{CDF of the delay of a iperf3 measurement.}
    \label{fig:eval:cdf.iperf.delay}
\end{minipage}
\begin{minipage}[t]{0.49\textwidth}
    \centering
    \includegraphics[width=0.99\textwidth]{/mnt/ext1/5g-masterarbeit-daten/main_measurement/.plots/CDF_delay_9ff091db__4b3d60f0_full_4:3.pdf}
    \caption{CDF of the delay of a scapy ping measurement.}
    \label{fig:eval:cdf.scapy.delay}
\end{minipage}%
\end{figure}

Throughout this work, achieved performance is compared on a per-run basis, aggregated to get a mean and a \gls{CI}
for this paramter combination.
While the throughput is a result of the transmitted data volume over the duration of a run, the average delay of a run is in itself a mean value.
In this section, the justification for this approach is presented.
In Figure~\ref{fig:eval:cdf.iperf.delay} and Figure~\ref{fig:eval:cdf.scapy.delay}, the \gls{CDF} of the per run delay is shown.
In Figure~\ref{fig:eval:cdf.iperf.delay}, the delay of an iperf3 measurement with maximum possible throughput can be seen,
while the delay of a ping measurement with small packets and an \gls{IAT} of \SI{10}{\milli\second} can be seen in
Figure~\ref{fig:eval:cdf.scapy.delay}.
In both figures, the different runs are highlighted in different colors, the delay is marked on the x-axis and the \gls{CDF} is marked on the y-axis.

The figures indicate, that the runs do not all follow the same distribution.
Variability in latency can be seen, which may be attributed to system noise or uncontrolled factors.

In generall, there are multiple parameter combinations, where all runs behave similarly.
However, for multiple cases there are clear differences in the distribution of delay times and runs show significant different metrics as a result.
We therefore decided, to not concatenate repetitions, but use a per-run average as a basis.

\subsection{UHD and gNB versions}\label{sec:eval:versions}

In this section, the influence of the driver and \gls{gNB} software version is investigated.
In Figure~\ref{fig:eval:uhd.version.throughput}, the throughput for two different \gls{UHD} versions is compared.
Vertical facets differentiate between uplink and downlink and horizontal facets, as well as the color is used to identify the \gls{TDD} pattern.
Throughput in \SI{}{\mega\bit\per\second} is marked on the y-axis and the \gls{UHD} versions are displayed on the x-axis.
There is no statistically significant improvement in throughput performance for a newer \gls{UHD} version, however, a positive trend can be seen for some uplink configurations.


\begin{figure}[htb]
\centering
\includegraphics[width=0.95\textwidth]{/mnt/ext1/5g-masterarbeit-daten/main_measurement/.plots/UHD__iperf_throughput_alt01.pdf}
\caption{Influence of different \gls{UHD} versions on the throughput performance.}
\label{fig:eval:uhd.version.throughput}
\end{figure}


The influence of the \gls{gNB} version on the throughput can be seen in Figure~\ref{fig:eval:gnb.version.throughput}.
Different \gls{gNB} versions are indicated on the x-axis and downlink and uplink traffic is differentiated by vertical facets.
Horizontal facets in combinations with the color coding are used to signal the \gls{TDD} pattern.
The throughput performance for OAI 2.2.0 is significantly improved compared to version 2.1.0.
All downlink configurations show statistically significant improvements.
Conversely, uplink performance is only significantly improved in one case: for a \gls{TDD} length of five slots and a ratio of 1:1.
For srsRAN, a change from version 24.04 to 24.10 results in no significant improvement in downlink performance.
Uplink throughput, on the other hand, is significantly improved for multiple configurations, in particular for a \gls{TDD} length of five slots.




\begin{figure}[htb]
\centering
\includegraphics[width=0.95\textwidth]{/mnt/ext1/5g-masterarbeit-daten/main_measurement/.plots/GNB__iperf_throughput.pdf}
\caption{Influence of different \gls{gNB} versions on the throughput performance.}
\label{fig:eval:gnb.version.throughput}
\end{figure}



\subsection{Effect of TDD Pattern}\label{sec:eval:tdd}

In Section~\ref{sec:back:frame} and Section~\ref{sec:back:duplexing}, the 5G frame structure and uplink downlink duplexing is explained.
In this work, duplexing is explained by two parameters, the \gls{TDD} period and the downlink-to-uplink ratio.
From a theoretical standpoint, the \gls{TDD} period is expected to have a minor effect on the throughput, since longer periods reduce the relative overhead introduced by guard symbols.
The downlink-to-uplink ratio, on the other hand, should have a significant effect on the throughput,
with a 4:1 configuration favoring downlink throughput at the cost of reduced uplink throughput performance.

In Figure~\ref{fig:eval:tdd.throughput}, the impact of the \gls{TDD} configuration on the throughput of the tested can be seen.
The \gls{TDD} pattern is indicated by the color coding and x-axis ticks.
Horizontal and vertical facets differentiate between the type of \gls{gNB} and the traffic direction, respectively.
The expected behavior is visible for downlink srsRAN measurements showing a proportional increase in throughput for higher ratios.
However, for a \gls{TDD} period of 20 slots, there is one outlier at a ratio of 4:1, showing significantly less performance then expected.
Uplink measurements for srsRAN display the expected behavior only in the case of a 10 slot \gls{TDD} period.
For uplink and OAI type measurements, the exected proportional decrease of achieved throughput can only be observed for a period length of five slots.
No definite statement can be made for a \gls{TDD} period of 10 slots, because this period can't be evaluated with OAI for a ratio of 1:1.
Downlink OAI measurements show the expected behavior for no \gls{TDD} period, instead, the best throughput can be achieved for a ratio of 2:1.


\begin{figure}[htb]
\centering
\includegraphics[width=0.95\textwidth]{/mnt/ext1/5g-masterarbeit-daten/main_measurement/.plots/TDD__iperf_throughput_alt01.pdf}
\caption{Influence of \gls{TDD} pattern on the throughput performance.}
\label{fig:eval:tdd.throughput}
\end{figure}


\subsection{Comparison of Measurement Results}\label{sec:eval:cnsm}

In Raffeck et al. \cite{raffeck2024parameterizing} the measurement results of three different testbeds are compared,
one testbed at the Norwegian University of Science and Technology and two testbeds at the University of Würzburg.
The testbeds located at the University of Würzburg differ by the height difference between the \gls{UE} and the \gls{RU}.
In one case, \gls{UE} and \gls{RU} are located at the same height and in on case, the \gls{RU} is elevated above the \gls{UE}.

Table~\ref{tab:eval:cnsm} presents a comparison between the measurement results obtained in this master's thesis (MT)
and those from two other testbeds: the same-height setup at the University of Würzburg (UWUE) and the testbed at the
Norwegian University of Science and Technology (NTNU).
The results from this work are selected to match the application versions used in the paper and are aggregated across all \gls{TDD} patterns.
The measurements conducted for this work differ significantly to the UWUE measurements in seven of 12 cases and four out of 12 cases when compared to the NTNU testbed.
Significant differences were observed in seven out of twelve cases when compared with UWUE and in four out of twelve cases relative to NTNU.
Despite these differences, all three testbeds exhibit consistent trends.
The most notable being that the OAI downlink performance peaks at a 2:1 downlink-to-uplink ratio.

\begin{table}[htb]
\centering
\caption{Measurement results compared to Raffeck et al.\cite{raffeck2024parameterizing} Green indicates no significant difference and orange indicates a significant difference to the measurements of this thesis.}
\label{tab:eval:cnsm}
\scriptsize

\begin{tabular}{c | c | c | c  c | c | c  c}
    \toprule
    \multirow{2}*{gNB}  & \multirow{2}*{Ratio} & \multicolumn{3}{c | }{ DL (Mbps) } & \multicolumn{3}{c}{ UL (Mbps) } \\
                        &                      & MT & UWUE & NTNU & MT & UWUE & NTNU \\
    \midrule
    \multirow{3}*{srsRAN}    & 1:1   & {{  44.07$\pm$1.31  }}    & {\setulcolor{Cgreen3} \ul{  42.87$\pm$1.04  }}  &  {\setulcolor{Cgreen3} \ul{  42.81$\pm$1.03  }}  & {{  19.97$\pm$4.68  }}     & {\setulcolor{Cgreen3} \ul{  18.05$\pm$2.23   }} & {\setulcolor{Corange3} \ul{ 8.96$\pm$3.90  }} \\
                             & 2:1   & {{  57.09$\pm$0.47  }}    & {\setulcolor{Corange3}\ul{  55.54$\pm$0.40  }}  &  {\setulcolor{Corange3}\ul{  55.64$\pm$0.43  }}  & {{  10.62$\pm$4.39  }}     & {\setulcolor{Corange3}\ul{  17.21$\pm$1.61   }} & {\setulcolor{Corange3}\ul{ 5.04$\pm$2.72  }} \\
                             & 4:1   & {{  61.67$\pm$4.55  }}    & {\setulcolor{Cgreen3} \ul{  60.02$\pm$3.52  }}  &  {\setulcolor{Cgreen3} \ul{  59.86$\pm$3.50  }}  & {{   8.90$\pm$2.28  }}     & {\setulcolor{Corange3}\ul{  13.98$\pm$0.10   }} & {\setulcolor{Cgreen3} \ul{ 5.94$\pm$3.14  }} \\
    \midrule
    \multirow{3}*{OAI}       & 1:1   & {{  46.78$\pm$1.49  }}    & {\setulcolor{Corange3}\ul{  38.37$\pm$1.83  }}  &  {\setulcolor{Cgreen3} \ul{  42.20$\pm$0.84  }}  & {{  10.04$\pm$4.21  }}     & {\setulcolor{Cgreen3} \ul{   9.51$\pm$3.37   }} & {\setulcolor{Cgreen3} \ul{11.31$\pm$4.42  }} \\
                             & 2:1   & {{  60.18$\pm$4.55  }}    & {\setulcolor{Corange3}\ul{  49.51$\pm$4.04  }}  &  {\setulcolor{Cgreen3} \ul{  55.10$\pm$4.21  }}  & {{   8.21$\pm$1.53  }}     & {\setulcolor{Cgreen3} \ul{   7.17$\pm$1.11   }} & {\setulcolor{Cgreen3}\ul{ 8.47$\pm$1.29  }} \\
                             & 4:1   & {{  42.54$\pm$0.10  }}    & {\setulcolor{Corange3}\ul{  35.87$\pm$0.80  }}  &  {\setulcolor{Corange3}\ul{  39.25$\pm$1.13  }}  & {{   8.17$\pm$0.24  }}     & {\setulcolor{Corange3}\ul{   7.13$\pm$0.10   }} & {\setulcolor{Cgreen3} \ul{ 8.25$\pm$0.07  }} \\
    \midrule
\end{tabular}
\end{table}


\subsection{ANOVA Results}\label{sec:eval:anova}

To further investigate the impact of parameters selected for the case study, two full factor \gls{ANOVA} are conducted,
one for throughput measurements and one for delay measurements.
The \gls{ANOVA} quantifies the proportion of the total variance in the measurement results that can be attributed to a certain parameter or parameter interaction.
The significance of these contributions is evaluated using an F-test.


\begin{table}[htb]
\centering
\caption{Shortened results of a full-factor \gls{ANOVA} for throughput measurements.}
\label{tab:eval:anova.iperf.throughput}
\begin{tabular}{l | c | c }

    \toprule
Param                                     & .p        & SS/SST \\
    \midrule
\texttt{direction}                                 & < 0.0001  & 0.8516 \\
\texttt{tdd\_ratio} : \texttt{direction}                       & < 0.0001  & 0.0317 \\
\texttt{gnb\_type} : \texttt{tdd\_ratio} : \texttt{direction}              & < 0.0001  & 0.0179 \\
\texttt{gnb\_type} : \texttt{tdd\_ratio}                        & < 0.0001  & 0.0141 \\
\texttt{gnb\_type}                                  & < 0.0001  & 0.0112 \\
\texttt{tdd\_period}                                & < 0.0001  & 0.0110 \\
\texttt{tdd\_ratio}                                 & < 0.0001  & 0.0095 \\
\texttt{tdd\_ratio} : \texttt{tdd\_period}                      & < 0.0001  & 0.0087 \\
\texttt{tdd\_ratio} : \texttt{tdd\_period} : \texttt{direction}            & < 0.0001  & 0.0086 \\
\texttt{gnb\_type} : \texttt{tdd\_period}                       & < 0.0001  & 0.0062 \\
\texttt{gnb\_version}                               & < 0.0001  & 0.0055 \\
\texttt{[...]} & & \\
Residual                                  &  {-}      & 0.0023 \\

\end{tabular}
\end{table}



The \gls{ANOVA} of the throughput measurements can be seen in Table~\ref{tab:eval:anova.iperf.throughput}.
This table presents a condensed summary, while Table~\ref{tab:apendix:anova.iperf.throughput} lists an extended report.
The first column lists the parameters and parameter interactions, the third column displays the proportion of total variation,
explained through this parameter or parameter combination, and the second column states the statistical significance of this attribution.

The most substantial portion of the variance can be attributed to the direction of the traffic.
Given the consistently higher downlink throughput observed in this work compared to uplink, this is expected.

In Section~\ref{sec:eval:tdd}, the influence of the \gls{TDD} ratio is discussed.
The second biggest factor for the total throughput variance is the interaction between the \gls{TDD} ratio and the traffic direction.
In theory, a higher ratio should lead to diverging behavior, increasing the downlink and simultaneously decreasing the uplink throughput.
While this behavior can not be observed reliably, the interaction between the ratio and the direction is still substantial.
Further, in Section~\ref{sec:eval:tdd}, the different behavior between OAI and srsRAN in regard to the \gls{TDD} ratio is presented.
This behavior is also reflected by the third biggest factor for the total throughput variance: the interaction
between the \gls{gNB} type, the \gls{TDD} ratio and the traffic direction.
The row listed as \emph{Residual} lists the remainder which can not be contributed to one parameter or parameter interaction.
It therefore quantifies the portion of total variance which is explained by variation between run repetitions.




\begin{figure}[htb]
\centering
\includegraphics[width=0.95\textwidth]{/home/lks/DocSync/Uni/5G-Masterarbeit/analysis/main_interactions_4_iperf_throughput_1.pdf}
\caption{Interaction plot of throughput for strongest 2-factor interactions.}
\label{fig:eval:anova.interaction.iperf.throughput}
\end{figure}

The four strongest 2-factor interactions, as highlighted by the \gls{ANOVA}, are visualized in an interaction plot in
Figure~\ref{fig:eval:anova.interaction.iperf.throughput}.
The two parameter types form the horizontal and vertical facets, while the parameter values are colored and marked on the
x-axis, respectively.
The throughput in \SI{}{\mega\bit\per\second} is shown on the y-axis.
The strongest two-factor interaction is between \texttt{tdd\_ratio} and \texttt{direction}, accounting for \SI{3.2}{\percent}
of the total variation.
As explained in Section~\ref{sec:eval:tdd}, in theory, there should be an even stronger effect visible.
A high downlink to uplink ratio should increase the downlink and reduce the uplink, since more slots are
reserved for downlink and less slots are reserved for uplink.
However, the observed effects are less pronounced than expected.
Notably, the highest downlink throughput was achieved with a $2:1$ ratio rather than a $4:1$ ratio.

Another noteable effect can be seen for the interaction between the \texttt{tdd\_ratio} and the \texttt{gnb\_type},
which accounts of \SI{1.4}{\percent} of the total variation.
The extent to which a configured \gls{TDD} ratio improves throughput depends on the selected \texttt{gnb\_type}.
Specifically, OAI shows a significant drop in throughput at a 4:1 \gls{TDD} ratio, whereas srsRAN maintains higher throughput performance under the same conditions.

Additionally, the effect of the \gls{TDD} period is highlighted again.
Theoreticaly, there should be no significant impact of the interaction between \gls{TDD} period and ratio, however a distinct interaction is visible.





\begin{table}[htb]
\centering
\caption{Shortened results of a full-factor ANOVA for delay measurements.}
\label{tab:eval:anova.scapy.delay}
\begin{tabular}{l | c | c }

    \toprule
Param                                     & .p        & SS/SST \\
    \midrule
\textit{Residual}                                                     & {-}       & 0.2512 \\
\texttt{direction}                                                    & < 0.0001    & 0.0397 \\
\texttt{size}                                                         & < 0.0001    & 0.0375 \\
\texttt{size} : \texttt{direction}                                    & < 0.0001    & 0.0370 \\
\texttt{iat}                                                          & < 0.0001    & 0.0356 \\
\texttt{iat} : \texttt{direction}                                     & < 0.0001    & 0.0347 \\
\texttt{size} : \texttt{iat}                                          & < 0.0001    & 0.0341 \\
\texttt{size} : \texttt{iat} : \texttt{direction}                     & < 0.0001    & 0.0302 \\
\texttt{gnb\_type} : \texttt{size}                                     & < 0.0001    & 0.0147 \\
\texttt{gnb\_type} : \texttt{size} : \texttt{direction}                & < 0.0001    & 0.0146 \\
\texttt{gnb\_type} : \texttt{direction}                                & < 0.0001    & 0.0141 \\
\texttt{gnb\_type} : \texttt{iat}                                      & < 0.0001    & 0.0140 \\
\texttt{gnb\_type} : \texttt{iat} : \texttt{direction}                 & < 0.0001    & 0.0137 \\
\texttt{gnb\_type}                                                     & < 0.0001    & 0.0136 \\
\texttt{[...]} & & \\

\end{tabular}
\end{table}




The \gls{ANOVA} results for the ping-type delay measurements is presented in Table~\ref{tab:eval:anova.scapy.delay}, shortened to the most influential factors.
A more extensive presentation is given in Table~\ref{tab:appendix:anova.scapy.delay}, which lists the 30 most influential factors.
In contrast to the throughput measurements, a substantial portion of the total variance in delay can be attributed to variation between repetitions.
Further, the most significant factors include traffic parameters, like direction, packet size, and \gls{IAT},
while \gls{gNB} type and \gls{TDD} configuration are secondary.

The two parameter types define the horizontal and vertical facets of the plot, while the parameter values
indicated by color and x-axis placement.
The delay in \SI{}{\second} is shown on the y-axis.

The most significant interaction is between the traffic direction and the packet size.
A bigger packet size significantly increases the delay in the uplink direction,
however no effect can be seen for the downlink direction.
This behavior can be seen in multiple cases.
The combination of two parameters which influence the system load provide an disproportionate impact on the delay.
These parameters are a higher packet size, a smaller \gls{IAT}, or uplink
instead of downlink traffic, since downlink provides better throughput performance.
Additionally, OAI is more susceptible to be impacted by one of the other parameters.



\begin{figure}[htb]
\centering
\includegraphics[width=0.95\textwidth]{/home/lks/DocSync/Uni/5G-Masterarbeit/analysis/main_interactions_4_scapy_delay_2.pdf}
\caption{Interaction plot of delay for strongest 2-factor interactions.}
\label{fig:eval:anova.interaction.scapy.ping}
\end{figure}





\chapter{Future Work}\label{chap:future}

While this thesis provides a comprehensive evaluation of the key influence factors of the reproducibility and
behavior of 5G campus deployments, several limitations apply.
These limitations are not only technical constraints encountered during the experimental process but also
serve as starting points for future investigations.

In the following, potential directions for future work are discussed.
First, challenges related to the time requirements of the measurements are outlined.
Subsequent sections explore how a broader set of hardware components on the \gls{gNB} and the \gls{UE} side,
could help isolate hardware-specific effects and improve the generalizability of the results.
Additionally, future research should address the limited reliability and resolution of channel quality indicators,
to provide more context and a better explainability of unexpected results.
Further avenues include a more fine-grained analysis of the performance tuning and a controlled channel degradation using a jammer.



\section{More Measurements}\label{sec:future:moremsm}
A significant hinderance throughout this work is the long reset time of the USB modems, resulting in nearly
\SI{4}{\minute} of required time for \SI{1}{\minute} of actual traffic.
As a result, the factorial design for the case study with the parameters \texttt{repetitions},
\texttt{\gls{TDD} pattern},
\texttt{\gls{gNB} x \gls{UHD} versions}, \texttt{traffic direction for iperf3}, and
\texttt{\gls{IAT} x packet sizes for scapy ping} results in
$ 10 \cdot 9 \cdot 8 \cdot (2 + 2\cdot2) \cdot 4\text{min} \overset{\scriptscriptstyle\wedge}{=}$ 12 days for
this measurement alone.
The development of the measurements methodology consumed a significant portion of the project timeline,
primarily due to the detection and mitigation of the issues described in Section~\ref{sec:meth:issues}.
Therefore, the inclusion of additional parameters was not feasable within the scope of this thesis.
Additionally, the included parameters for the different scenarios were picked carefully with the time requirement of
the measurements in mind.
In the future, more measurements with an expanded set of controlled parameters should be performed.


\section{Different Radio Units}\label{sec:future:modem}
A single \gls{SDR} has been utilized in this work.
In the future, comparisons to different \gls{SDR}, as well as different versions of the same \gls{SDR},
should be performed to quantize the effect of the relevant hardware components.
In addition, exploring a proprietary \gls{gNB} and contrasting it with the open-source stacks used here could offer
deeper insights into the role of configuration and vendor-specific optimizations on system behavior.


\section{Different Modems}\label{sec:future:modem}

In this work, the SIM8200EA-M2 5G HAT modem was deployed.
Additionally, Quectel RM500Q and Quectel RM520N based modems were tested.
However, no connection between the modem and the \gls{gNB} could be established.
From the MAC layer PCAP recordings, it was visible, that no \gls{RRC} messages where received by the \gls{gNB}.
This problem could be replicated for a different testbed, not part of this work.
Given the time constraint, further debugging could not be justified.
In future, however, the influence of different modems regarding 5G \gls{NR} should be addressed.


\section{Channel Quality Indicators}\label{sec:future:channelquality}
The collection and evaluation of channel quality indicators, alongside the \gls{QoS} metrics, serves as an important
part in creating reproducible experiments and consistent results in 5G research.
However, three significant limitations regarding the channel metric acquisition were identified in this work.
First, as outlined in Section~\ref{sec:meth:issueqmi}, the modem's channel metrics could not be collected reliably,
limiting the ability to correlate \gls{UE}-side observations with network behavior.
Second, as discussed in Section~\ref{sec:eval:gain}, the \gls{SNR} values reported by OAI appear implausible,
casting doubt on the overall reliability of the provided channel metrics by \gls{OAI}.
Third, both srsRAN and \gls{OAI} use a relatively large reporting intervall of one second.
This limits the detection and analysis of short-term effects in the radio link.
Addressing these limitations should be a priority in future work, to improve the interpretability and correctness of the measurements.


\section{Performance Tuning}\label{sec:future:perfromance}
Performance Tuning is regarded as a binary parameter for the evaluation.
However, it is comprised of multiple different steps, which are explained in Section~\ref{sec:meth:parameters}.
In the future, a more fine-grained approach could be used to identify the relevant configurations of the
\gls{gNB} host system and provide a more complete understanding of the performance and hardware requirements of
open-source 5G campus deployments.



\section{Jammer}\label{sec:future:jammer}
For this work, the influence of a jammer on the air interface was briefly observed.
By utilizing a second \gls{SDR} to send random noise on the same frequency band the \gls{RAN} is using, the channel quality can be impaired,
which in turn causes the \gls{gNB} to select a lower \gls{MCS} scaling, in order to compensate.
Finding a suitable gain setting of the jammer, however, proved difficult.
The gain range which negatively affect the transmission, however still allows for successfull transmissions is relatively small.
A gain value to low does not affect the transmission in measureable way.
A gain value to high results in a complete connection loss, either by disconecting the \gls{UE}, in case of srsRAN, or by
causing a segmentation fault, in case of OAI.
Further, the gain range which negatively affects the transmission, while still allowing transmissions, depends,
among others, on the type of \gls{gNB}, the traffic direction, the location of the jammer, and even the traffic pattern.
Therefore, no parameterized, automated setup of the jammer is included in the evaluation.
In the future, however, the ability to deliberately \todo{jam} the transmission is a significant aspect to
simulate sub-optimal channel conditions and could provide valuable
insights into the system behavior of small-scale 5G deployments.
